\documentclass[14pt,a4paper]{extarticle}
\usepackage[left=3cm,right=2cm,top=2cm,bottom=2cm,headheight=18pt]{geometry}
\usepackage{fontspec}
\setmainfont{Liberation Serif}
\setsansfont{Liberation Serif}
\setmonofont{Liberation Mono}
\usepackage{unicode-math}
\setmathfont{texgyretermes-math.otf}
\usepackage{amsmath,graphicx,booktabs,longtable,array,enumitem}
\usepackage[normalem]{ulem}
\usepackage{titlesec,fancyhdr}
\usepackage[hidelinks,unicode]{hyperref}
\usepackage{url}
\urlstyle{same}
\setlength{\parindent}{1.27cm}
\setlength{\parskip}{0pt}
\linespread{1}
\setlength{\emergencystretch}{2em}
\clubpenalty=10000
\widowpenalty=10000
\titleformat{\section}{\normalfont\bfseries}{\thesection.}{0.5em}{}
\titleformat{\subsection}{\normalfont\bfseries}{\thesubsection.}{0.5em}{}
\titlespacing*{\section}{0pt}{1.2em}{0.5em}
\titlespacing*{\subsection}{0pt}{0.9em}{0.3em}
\setlist{nosep,leftmargin=1.27cm}
\pagestyle{fancy}
\fancyhf{}
\renewcommand{\headrulewidth}{0pt}
\renewcommand{\footrulewidth}{0pt}
\fancyhead[C]{\ifnum\value{page}>1\thepage\fi}
% Liberation Serif is a disclosed Times-compatible substitute, not Times New Roman.
% Current statutory geometry/size baseline: O‘RQ-682 appendix §§43–51.
\hyphenpenalty=10000
\exhyphenpenalty=10000
\pretolerance=10000
\tolerance=10000
\setlength{\emergencystretch}{4em}
\renewcommand{\normalsize}{\fontsize{14bp}{16.8bp}\selectfont}
\normalsize
\begin{document}
\begin{center}\textbf{TUSHUNTIRISH XATI}\end{center}
Asosiy ruscha murojaatga mos o‘zbekcha tarjima. 2-material; 1-material normativ qoidalar, 3-material ilmiy asos va elektron hisoblash qo‘shimchasidir.
\noindent Ijtimoiy reyestr hisob-kitoblarini tekshirish tartibi bo‘yicha mualliflik taklifiga ilmiy-texnik asos. Tahrir sanasi: 2026-yil 1-oktabr; hisoblash tayanch sanasi: 2026-yil 30-sentabr.\par
\section*{Tayyorlash asosi, ishlab chiquvchilar va maqsad}
PF-193-son Farmonning 13–14-bandlari 2027-yil 1-yanvardan ko‘p o‘lchovli baholashni va tegishli loyihani 2026-yil 1-dekabrgacha Prezident Administratsiyasiga kiritishni nazarda tutadi. Rasmiy ishlab chiquvchilar — O‘zbekiston Respublikasi Prezidenti huzuridagi Ijtimoiy himoya milliy agentligi va O‘zbekiston Respublikasi Iqtisodiyot va moliya vazirligi. Ushbu tushuntirish xati Abduxoliq Akmal ugli Ashuraliyevning mualliflik taklifiga tegishli; u rasmiy ishlab chiquvchi nomidan taqdim etilgan xat emas.

Maqsad — tasdiqlangan baholash qoidasini bir xil kirish qiymatlaridan qayta hisoblash, chegaralar va istisnolarni tekshirish, ma’lumotdagi muammoni dasturiy xatodan ajratish hamda individual qaror asoslarini tushuntirish. Asosiy taklif yetti asosiy band va yigirma sakkiz bandli ilovadan iborat. U yangi 2027-yilgi ballar yoki yordam miqdorlarini o‘zi belgilamaydi. Kutilayotgan natija — tekshiriladigan va qayta tiklanadigan hisob-kitob; oilalar farovonligi, budjet tejalishi yoki yordam qamrovi uchun miqdoriy prognoz berilmagan. To‘liq shakl o‘rniga qo‘shib qo‘llanmaydigan besh bandli qisqa muqobil ham berilgan.

Huquqiy shakl asosiy loyiha doirasida belgilanadi. PF-258-son Farmon mezonlariga tegadigan variant Prezident hujjatiga o‘zgartirishni talab qiladi. Vazirlar Mahkamasi hujjati yoki dasturiy sozlama yuqori yuridik kuchdagi mezonni almashtira olmaydi.
\section*{Tartibga solish predmeti va taklif etilgan hujjat tuzilishi}
Asosiy yechim — yangi ehtiyoj modelini tanlash emas, tasdiqlangan qoidalar bajarilishini majburiy tekshirish. Yetti asosiy band tartibni tasdiqlash, yagona tavsif, qayta bajarish, 22-modda bo‘yicha tayyorlash, bog‘liq hujjatlarni muvofiqlashtirish, resurs hisobi va qo‘llashdan oldingi tayyorlikni belgilaydi. Ilova yetti bobdan iborat: umumiy qoidalar (1–3-bandlar); tavsif va talqinlar (4–7); ma’lumot sifati va uzilishlar (8–11); hisoblashni tekshirish (12–17); qaror asoslari va qayta ko‘rish (18–20); tadqiqot, hisobot va nazorat (21–24); qabul, tuzatish va saqlash (25–28). 1-material normativ talablarni, ushbu xat maqsad va ijroni, 3-material batafsil dalillarni beradi.

Taklif PF-193 topshirig‘ini qonunosti tartibida bajarishga taalluqli. Yagona uslubiyot 84-bandidagi qonunosti munosabatlarini qonun darajasiga ko‘tarishni ko‘rib chiqish talabi qonun loyihasini tayyorlashga tegishli; bunda alohida qonun loyihasi taklif qilinmaydi. Asosiy hujjat turini vakolatli organ tanlaydi. Yordamga yangi huquq yoki to‘lov miqdori tegishli huquqiy kuchdagi hujjat bilan belgilanadi. Hisoblash tekshiruvi shu vakolatni almashtirmaydi.

\section*{Har bir talablar guruhining zarurati}
35-son qaror 1-ilovasi 46-bandidagi to‘g‘ri ishlash va o‘z vaqtida ma’lumot berish vazifalari boshlang‘ich asosdir. Tartib bajarishni tekshirishga imkon beradi: muayyan talqinning qo‘llangan qoidasi, kirish asosi, oraliq solishtirish va ziddiyatni hal qilgan shaxs aniqlanadi. Shu banddagi arizachi, ijtimoiy xodim va “mahalla yettiligi” vazifalari saqlanadi; davlat manbasining javobgarligi oilaga o‘tkazilmaydi.

Bir xil son turli birlik, davr va kiritish shartida bir xil kirish bo‘lmagani uchun mezon pasporti va sanali talqin zarur. Sifat holatini ajratish olinmagan yozuvni tasdiqlangan nol daromadga yoki keyinroq olishni ko‘proq ishonchlilikka aylantirmaydi. Istisno va ustuvorlik nazorati to‘g‘ri umumiy formula maxsus shart o‘tkazib yuborilgani sabab noto‘g‘ri qaror berishini oldini oladi. Komponent solishtirish jami natijada bir-birini qoplagan xatoni topadi. Individual izoh va bog‘langan tuzatish faqat statistika emas, muayyan qaror asosi va qonuniy qayta ko‘rishni ta’minlaydi.

Tartib zarur eng kam ma’lumot va amallarni belgilaydi. Agentlik mezon, murojaat muddati yoki rad asosini o‘zgartirmasdan ichki pasport, bayonnoma va jurnal shaklini belgilashi mumkin. Ishlab chiquvchi teng amaldagi talab va ishlayotgan nazoratni ko‘rsatsa, band unga qo‘shiladi. Tizim yoki umumiy nazoratga havola birlik, chegara, istisno va izohdagi alohida isbotlangan masalani hal etmaydi.

\section*{Har bir band uchun huquqiy vositani tanlash}
O‘RQ-682ning 20-moddasi mavjud mexanizm va ma’muriy vakolatlar yetarliligini tekshirishni talab qiladi. PF-193 topshirig‘i loyiha vazifasini belgilaydi, har bir nazoratga yangi norma kerakligini isbotlamaydi. Jadvalda N — vakolatli normativ qabul talab qiladigan taklif etilgan ommaviy majburiyat yoki vakolat; M — tengligi isbotlansa takrorlanmasdan birlashtiriladigan mavjud vazifa; I — mavjud vakolatdagi ichki ijro. Aralash qayd ommaviy minimum va ijroni ajratadi. Bu muallifning asoslangan tasnifi, Agentlik vakolatsizligining rasmiy xulosasi emas. 35-son qaror 1-ilovasining 45–50-bandlari va shaxsga doir ma’lumotlar Qonunining 11 va 24-moddalari mavjud vositalardir.

\begingroup\fontsize{11bp}{13.2bp}\selectfont
\begin{longtable}{p{0.12\linewidth} p{0.13\linewidth} p{\dimexpr0.75\linewidth-6\tabcolsep\relax}}
\toprule Band & Tanlov & Asos va yechim\\\midrule\endhead
Asosiy 1 & N/I & Ommaviy minimumni qabul qilish, shakllarni topshirish; mavjud nazorat nomini o‘zgartirish alohida hujjatni talab qilmaydi.\\
Asosiy 2 & N/M & Mezon vakolatli tasdiq talab qiladi; tavsif uni o‘zi o‘zgartirmaydi.\\
Asosiy 3 & M/I & Tizimning to‘g‘ri ishlashi va muddatlar mavjud; sinov mavjud vakolatda.\\
Asosiy 4 & M/I & 22-modda va ichki reglament; yangi tashqi malaka yoki fuqaro vazifasi yo‘q.\\
Asosiy 5 & M/I & Qoida/tartib muvofiqligi; zaxira kiritish yoki boshlanish o‘zgarishi yo‘q.\\
Asosiy 6 & M/I & Qo‘shimcha vazifa, mavjud qamrov va haqiqiy mablag‘; takroriy xarid yo‘q.\\
Asosiy 7 & M/I & Chiqarilish qaydi va kelishuv; yangi hisobot ishga tushish sharti emas.\\
Ilova 1 & I & Qamrov va bog‘langan nazorat mavjud mas’uliyatni tashkil etadi.\\
2 & I & Atamalar bir xil ijro uchun, mustaqil yordam huquqi emas.\\
3 & M & Toifa va yordam shartlari alohida; soha qoidalari saqlanadi.\\
4 & I & Pasport va manba tafsiloti ichki reglamentga o‘tadi.\\
5 & N/I & Ommaviy mezon tenglik va ustuvorlikni hal qiladi; dastur ijroni tekshiradi.\\
6 & I/M & Talqin va huquqiy sana mavjud vazifalar asosida tiklashni ta’minlaydi.\\
7 & M/I & Tasdiqlovchi akt va tavsif mavjud resursda; ichki talqin.\\
8 & I/M & Manba tarixi qonuniy asos va aniqlik uchun; qayd sxemasi ichki.\\
9 & N/M & Yo‘q ma’lumot yangi salbiy asos emas; qonuniy me’yoriy qiymat saqlanadi.\\
10 & N/M/I & Vakolatli aniqlashtirish, muddat va yangi rad etish yo‘qligi; yo‘llash shakli ichki.\\
11 & M/I & Mavjud uzilish vakolatini qayd qilish, sohasini kengaytirmaslik.\\
12 & I & Mustaqil nazorat, qamrov va tarkibni solishtirish to‘g‘ri ishlashni bajaradi.\\
13 & I & Oldingi reja natija talqinini boshqaradi, arizachi vazifasini emas.\\
14 & I & Surat va maxraj hisobot sifati uchun, muhtojlik mezoni emas.\\
15 & I & Muhtojlik aniqligi dalil talab qiladi; empirik tadqiqot ishga tushish sharti emas.\\
16 & I & Bayonnoma dalili va shakli ichki reglamentda.\\
17 & M/I & Nuqson va ta’sirlangan qarorlar; o‘zgarish qonuniy vakolatga bog‘liq.\\
18 & M/I & Ochiq qaror asoslari saqlanadi; texnik tarix shakli ichki.\\
19 & N/M & Mavjud izoh huquqi birlashtiriladi; hisob izohi SMSni to‘ldiradi.\\
20 & M/I & Qonuniy e’tiroz va tuzatish muddatlari; majburiy sudgacha bosqich yo‘q.\\
21 & M/I & Qonuniy ishlov va maxsus shartlar; tadqiqot ruxsati ichki.\\
22 & M/I & Himoyalangan ijro va oshkor tekshiruvi; tashqi shaxsiy ma’lumot kirishi yo‘q.\\
23 & M/I & Ichki bayonnoma va qonuniy kirish/oshkor; yangi davriy hisobot yo‘q.\\
24 & M & Mavjud nazorat; yangi organ yoki kamaytirish asosi yo‘q.\\
25 & I/M & Mavjud chiqarilish qaydi; tashqi majburiy xulosa saqlanadi.\\
26 & M/I & Dasturni amaldagi qoidaga tuzatish/sinash; qo‘lda zaxira yoki eskirgan mezon yo‘q.\\
27 & M & Sana va qonuniy qayta ko‘rish; o‘tish qarzidan ozod qilish yo‘q.\\
28 & M/I & Qonuniy saqlash va yo‘q qilish; toifa jadvali ichki.\\
\bottomrule\end{longtable}\endgroup

\section*{Normativ mazmun zarurati va texnik yechim chegarasi}
\textbf{Zarurat predmeti.} \href{https://lex.uz/docs/-8472526}{PF-193}ning 13–14-bandlari yangi ko‘p o‘lchovli ball bo‘yicha toifalashni belgilab, normativ loyihani topshiradi. Bu 2027-yil normativ mazmuni ustida ishlashning bevosita asosi, ammo alohida sinov hujjati yoki muallifning barcha 28 bandi muqarrarligini isbotlamaydi. Zarur mazmun allaqachon topshirilgan hujjatga kiritiladi; 1-materialdagi besh bandli qisqa muqobil shu shaklni beradi. Bu yerda 2027-yil ballari, vaznlari yoki to‘liq parametrlar o‘ylab topilmagan; ishlab chiquvchi ularni qonuniy belgilashi kerak.

\textbf{Texnik yechim nega huquqiy yechimni almashtirmaydi.} \href{https://lex.uz/docs/-5378966}{O‘RQ-682}ning 3, 14, 17, 18 va 37–40-moddalari normativ mazmun, vakolat, yuridik kuch, shakl, tegishli holatdagi ro‘yxat va rasmiy e’lonni ajratadi. Chegara, tenglik, davr, ruxsat etilgan kirish yoki normativ hisoblashning tanlovi o‘zi kiritish yoki rad etish shartini belgilasa, muvaffaqiyatli sinovning o‘zi uni asoslamaydi. Sinov tanlangan qoida bajarilganini tekshiradi; vakolatli hujjat shu qoidani oilaga qo‘llashning huquqiy asosini beradi. Bu ko‘rsatilgan normalardan muallif xulosasi, har bir dastur parametriga alohida normativ tasdiq kerakligi da’vosi emas. Qonuniy mavjud vakolat va belgilangan shartning texnik bajarilishi saqlanadi.

\textbf{Hujjatlarni ajratish.} \href{https://lex.uz/docs/-7766987}{2025-yil 9-oktabrda 3686-son bilan ro‘yxatga olingan Qoidalar}ning 94–95-bandlari huquq va kafolatlarga ta’sir qiluvchi, idoralararo yoki qabul qilgan organ tizimidan tashqaridagi tashkilotlar uchun majburiy idoraviy normativ hujjatlarni qamraydi. 96(b)-band yangi huquqiy normasiz oldingi tartib ijrosiga oid hujjatni ro‘yxatdan istisno qiladi; 96(v)–(j) texnik istisnolarni, jumladan boshqa hujjatni Adliya vazirligi tomonidan texnik toifaga kiritishni belgilaydi. “Texnik talab” nomi mazmunning huquqiy tabiatini o‘zi hal qilmaydi. Aksincha, haqiqiy texnik hujjat ushbu taklif tufayli normativga aylanmaydi. Qoidalar idoraviy shaklga taalluqli; Prezident yoki Vazirlar Mahkamasi hujjatini davlat ro‘yxatidan o‘tkazish talabi deb olinmaydi. 2026-yil 17-sentabr o‘zgarishli tahrir tekshirildi; 20-band zarurat va muqobillarni o‘rganishni bevosita talab qiladi.

\textbf{Tenglikning ikkita mustaqil sharti.} Huquqiy qamrov aniq band, vakolatli organ, yuridik kuch, oilalar/holatlar sohasi, qo‘llash sanasi va majburiy oqibatni, idoraviy shakl uchun tegishli ro‘yxat hamda e’lonni ko‘rsatadi. Texnik qamrov kirish, sana, istisno, mustaqil asoslangan javob, sinalgan talqin va bajarilgan natijani ko‘rsatadi. Barcha muallif sinovlari qamralsa, ularni takrorlash kerak emas. Qo‘shimcha norma faqat huquqiy mazmun ham to‘liq ta’minlanganda kerak bo‘lmaydi. Umumiy shartnoma, audit, guruh tarkibi buyrug‘i yoki qabul bayonnomasi yo‘q mezonning asosi deb olinmaydi. Ikkala qamrov isbotlansa, shu masala uchun qo‘shimcha norma taklif etilmaydi.

\textbf{Qo‘llanish sohasi va sharh.} Mavjud shartning vakolatli sharhi yoki manba ma’lumoti sohasining hujjati muayyan qarshi misolni hal qilishi mumkin. Manba, vakolat, davr, birlik va maydonning haqiqiy ma’nosi ko‘rsatiladi; matematik xulosa shu sohaga cheklanadi. Dasturning qiymatni hozir rad etishi oilaning tegishli faktni huquqan bildira olmasligini o‘zi isbotlamaydi. Yangi kiritish cheklovi imkonsiz kirish dalili sifatida joriy etilmaydi. O‘RQ-682ning 55-moddasi sharh orqali hujjatni o‘zgartirish, to‘ldirish yoki tuzatishga yo‘l qo‘ymaydi. Aniqlangan ma’no izchil ijroni, yangi huquqiy shart tanlovi vakolatli o‘zgartirishni talab qiladi. Yopilgan masala yangi norma uchun sun’iy asos sifatida saqlanmaydi.

\textbf{Bevosita mavjud vaqtinchalik vakolat.} \href{https://lex.uz/docs/-8027513}{35-son qaror} 1-ilovasi 48-bandi texnik nosozlik, uzilish va elektron ma’lumot to‘liq olinmaganini qamrab oladi. Agentlikka alohida tizim almashishini to‘liq tiklanishgacha vaqtincha to‘xtatish va toifani avtomatik aniqlash algoritmini o‘zgartirish huquqi beriladi. Shu sababli har vaqtinchalik algoritm o‘zgarishiga yangi hujjat kerak deb da’vo qilinmaydi. Qisqa 4-band va to‘liq 11-band vakolatni saqlab, muayyan qo‘llanishining kuzatiladigan qaydini taklif qiladi. Yangi uslubiyotga tatbiq doirasini ishlab chiquvchi muvofiqlashtiradi; 48-band 2027-yil butun doimiy ball tizimini belgilagan deb olinmaydi.

\textbf{Xarajat va tanlash ketma-ketligi.} Topshirilgan loyihada qonuniy baholash sharti va saqlanadigan kafolat aniqlanadi, bir vaqtda mavjud sinov va shartnomalar qiyoslanadi. So‘ng faqat qoplanmagan integratsiya baholanadi. Qisqa tahrir yangi ariza oqimi yoki respublika qo‘lda xizmatini yaratmaydi. Huquqiy tayyorlash, sinov bajarish va interfeys o‘zgarishi baribir vaqt talab qiladi; alohida mablag‘ yo‘qligi yoki ichki ijro va’dasi nol xarajatni isbotlamaydi. Haqiqiy soatlar, shartnoma qamrovi, quvvat va moliyalashtirish tasdiqlanmagan. Asoslangan resurs qiyinligi texnik tafsilot va ish taqsimotini o‘zgartirishi mumkin, lekin dasturga yo‘q yordam shartini belgilash huquqini bermaydi.

\textbf{Qabul mezoni.} Boshqa tartib, qisqaroq matn, amaldagi hujjatga havola yoki tegishli vakolatli idoraviy normativ hujjat qabul qilinadigan shakldir. Har huquqiy vazifa uchun bir xil majburiy oqibat, manba, sana va vakolat; texnik vazifa uchun ijrochi va teng dalil solishtiriladi. Bu muhokama uchun so‘ralgan qiyos, davlat javobining yangi majburiy shakli emas. Vakolatli e’lon qilingan hujjatdagi teng norma alohida ilovasiz ham normativ natija. Faqat texnik talabga misol kiritish texnik natija; u isbotlangan hal qilinmagan normativ mazmunni yopmaydi.

\begingroup\fontsize{11bp}{13.2bp}\selectfont
\begin{longtable}{p{0.26\linewidth} p{0.39\linewidth} p{\dimexpr0.35\linewidth-6\tabcolsep\relax}}
\toprule E’tiroz & Yetarli tasdiq & Taklif bo‘yicha qaror\\\midrule\endhead
Mavjud nazorat misolni qamraydi & Shu soha va sanaga tatbiq norma hamda mustaqil asoslangan bajarilgan sinov & Takror sinov olinmaydi; norma faqat huquqiy qamrov tasdiqlansa birlashtiriladi.\\
Qarshi misol soha tashqarisida yoki hal qilingan & Vakolatli cheklov yoki sharh, amal sanasi va kirish ma’nosi & Misol cheklanadi yoki olinadi; yangi huquqiy cheklov vakolatli hujjat talab qiladi.\\
Qo‘shimcha ish & Aniq qoplanmagan vazifa, soat, shartnoma, mablag‘ va ishlash o‘zgarsa tegishli yuklama & Texnik integratsiya mutanosib qilinadi; huquqiy shart sozlama bilan almashtirilmaydi.\\
Boshqa tahrir kerak & Shu vakolat, majburiy oqibat, soha, sana va ishlaydigan kanal & Qisqa yoki teng normativ tahrir natija deb qabul qilinadi; muallif matnini aynan olishning o‘zi maqsad emas.\\
\bottomrule\end{longtable}\endgroup

\textbf{Minimal qabul yo‘li.} Topshirilgan hujjatga zarur huquqiy shart va baholashga xos kafolat kiritiladi; 1-material besh bandli qisqa muqobilni beradi. To‘liq yetti band va ilova boshqa shakl, qo‘shiladigan talab emas. Qabul sinovi va batafsil shakl mavjud vakolat va teng qaydni ishlatadi. Huquqiy qamrov alohida tasdiqlansa, ilmiy misolning texnik kiritilishi muhandislik vazifasi uchun yetarli. U yo‘q huquqiy shartni qabul qilishni almashtirmaydi. Vakolat, majburiy oqibat, soha, sana, kirish va qonuniy muddat saqlansa qisqartirish yoki havola mumkin.

\section*{Dalil va u qo‘llab-quvvatlaydigan tor chora}
\begingroup\fontsize{11bp}{13.2bp}\selectfont
\begin{longtable}{p{0.30\linewidth} p{0.35\linewidth} p{\dimexpr0.35\linewidth-6\tabcolsep\relax}}
\toprule Bajarilgan dalil & Asoslangan tor chora & Isbotlanmagan narsa\\\midrule\endhead
Daromad tengligi va og‘ir vaziyat nazoratlari & Tenglik, istisno ustuvorligi va qonuniy kutilgan javobni muvofiqlashtirish. & Ishlayotgan tizim xatolari tezligi yoki qamrab oluvchi chegara siyosatining ustunligi.\\
4\,105 tuzilgan satrda 620 manfiy maydon natijasi & Vakolatli geometriya va sohani tasdiqlash, dasturlashdan oldin formulani hal etish va tanlangan qoidani tekshirish. & 620 zarar ko‘rgan oila, milliy chastota yoki nol chegarasining sof foydasi.\\
Rasmiy misol: 432\,000 so‘m & Birlik, oila bo‘luvchisi va davrni huquqiy misolga tekshirish. & Amaldagi tizimning umumiy to‘g‘riligi.\\
Soliq, chorva va shartli pul o‘tkazmasi nazoratlari & Tarmoqlar, birlik va sanali parametrni nazoratda saqlash. & Agentlik amalga oshirishi, to‘liq 2027 spetsifikatsiya yoki haqiqiy ma’lumot sifati.\\
37 regressiya sinovi va qat’iy tarkib tekshiruvi & Tuzilgan holatni tushirish yoki takrorlashdan himoyalangan qayta bajariladigan vosita. & Ishlab chiqarish tizimi qabuli, bo‘sh xodim yoki tejalgan mablag‘.\\
\bottomrule\end{longtable}\endgroup
Nazorat hal etish savolini belgilaydi, Agentlik ijrosiga nisbatan o‘lchangan qo‘shimcha foydani emas. Spetsifikatsiya, mustaqil javob, chiqarilish qaydi, izoh va zarur texnik/himoya xulosasi solishtiriladi. Teng nazorat qayta ishlatiladi; faqat aniqlangan qoplanmagan vazifa qo‘shimcha. Oila ma’lumoti Agentlikdan chiqmaydi. Vakolatli soha tuzilgan kirishni chiqarishi mumkin. Farovonlik bahosi alohida qamrovda. Yangi ochiq statistik hisobot taklif etilmaydi.

\section*{Normativ ilovadan tashqaridagi ichki texnik spetsifikatsiya}
Bu asosiy 4-band ichki reglamenti uchun muallif topshirig‘i, qo‘shimcha ommaviy mezon, arizachi vazifasi yoki qonunni o‘zgartirish vakolati emas.

\textbf{Mezon va manba pasportlari.} Barqaror identifikator, huquqiy qoida, tur va soha, egasi va yozuv identifikatori, birlik, zarurlik, davr chegarasi, yangilash, o‘zgartirish, takror va kesishma, ziddiyatli tanlash, qonuniy me’yoriy qiymat, amallar tartibi va oraliq aniqlik qayd etiladi. O‘tkazma va chegara bir birlik va davrda; oila tarkibi va o‘zgaruvchi ish haqi parametrida tegishli sana bo‘ladi. Qabul vaqti fakt sanasi emas. Ziddiyatli ikkala qiymat va qonuniy tanlov asosi saqlanadi. Xodim tuzatishida muallif, vaqt va vakolat qayd etiladi.

\textbf{Nazorat ro‘yxati.} Har bir qo‘llanadigan shart, chegara, istisno, ustuvorlik va erishiladigan o‘zaro ta’sir misolga bog‘lanadi. Misol identifikatori, kirishlar, dastlabki shartlar, amal sana, huquqiy chiqarilgan oraliq va yakuniy javob, chiqarish asosi va mustaqil tekshiruvchi saqlanadi. Rejalashtirilgan va bajarilgan identifikatorlar hamda takror sonlari solishtiriladi; chiqarish va qo‘llanmaslik izohlanadi. Bo‘sh tarkib va aniqlanmagan maxraj ochiq ko‘rsatiladi. Qoida yoki manba o‘zgarishi ta’sirlangan va bog‘liq yo‘llarni tekshiradi; soha aniqlanmasa butun qo‘llanadigan zanjir tekshiriladi. Dasturlarning umumiy javobi mustaqil huquqiy javob emas.

\textbf{Bayonnoma va tuzatish.} Uslubiyot, dastur va sozlama alohida identifikatsiyasi, ma’lumot davri va tarkibi, qamrov va ro‘yxat, kutilgan/haqiqiy komponentlar, bajarilgan/chiqarilgan sonlar, farq, ta’sir, tuzatish mas’uli va qayta sinov saqlanadi. Qoidaga moslik, ma’lumot sifati va qaror asoslari to‘liqligi ajratiladi. Nuqson faqat uni ko‘rsatgan shikoyatga emas, talqin, sana va arizalarga bog‘lanadi. Aniqlashtirish dastlabki qabul, manba yoki oluvchi, qonuniy muddat, javob, tuzatilgan hisob va bildirishni bog‘laydi. Himoyalangan yozuv ruxsat etilgan saqlashda qoladi; ommaviy qism faqat oshkor tekshiruvidan o‘tgan jamlanma beradi.

\section*{Yordamga yangi o‘tish huquqisiz ishga tushish muvofiqligi}
26-band endi tuzatish, qayta sinash va qonuniy dasturiy chiqarilishni boshqarishga tegishli. Ikkinchi qo‘lda qaror xizmati talab etilmaydi. Oldingi chiqarilish faqat qarorga qonuniy qo‘llanadigan qoidani bajarib, tegishli nazorat va ishlash yaroqliligi tekshirilsa ishlatiladi. Almashtirilgan daromad qoidasidagi chiqarilish eski testlardan o‘tgani uchun qonuniy zaxira emas. Uzilish chorasi 35-son qaror 1-ilova 48-band bilan cheklangan; boshqa nuqson bu vakolatni kengaytirmaydi.

Taqdim etiladigan normativ matnda 2026-yil uslubiyotini saqlash, daromad tengligini kiritish yoki o‘tish qarzidan ozod qilish bo‘yicha Prezident zaxira o‘tishi yo‘q. Bu tanlovlar huquq va budjet oqibatining alohida dalilini talab qiladi va tor hisoblash muvofiqligi hissasi uchun zarur emas. Quyidagi tenglik va yer formulasining quyi chegarasi variantlari shartli tadqiqot variantlari bo‘lib, qabul qilish sharti sifatida so‘ralmaydi. Zaxirani chiqarish dastur yo‘q normani to‘ldiradi degani emas: ishlab chiquvchi qonuniy qaror jadvalini tugatib, qoida yo‘qligi yoki ziddiyatini birinchi qo‘llanishdan oldin hal qiladi. O‘RQ-682 55-modda bo‘yicha vakolatli rasmiy sharh qonunosti normasini izohlaydi, o‘zgartirish kiritmaydi. Mezon yoki boshlanishni faqat vakolatli qabul qiluvchi o‘zgartiradi. Tayinlangan yordam, kutayotgan ariza va undirish amaldagi qonunchilik bilan tartibga solinadi; bu taklif oraliq huquq yaratmaydi.

\section*{Mavjud nazorat, amaliy dalil va qo‘shimcha vazifa}
Quyidagi rasmiy materiallar 2026-yil 1-oktabrda ochilib tekshirildi. Qonuniy vazifa xizmat yoki tadbir ishlashining dalilidan ajratiladi. Materiallarda Agentlikning to‘liq dasturiy chiqarilish sinov ro‘yxati yoki o‘lchangan resurs hisobi yo‘q.
\begingroup\fontsize{11bp}{13.2bp}\selectfont
\begin{longtable}{p{0.30\linewidth} p{0.33\linewidth} p{\dimexpr0.37\linewidth-6\tabcolsep\relax}}
\toprule Rasmiy dalil & Aniqlangan soha & Tor qo‘shish va chegara\\\midrule\endhead
35-son qaror, 1-ilova, 45–50-band & Himoyalangan ma’lumot, tizim/manba mas’uli, cheklangan uzilish vakolati, nazorat va nizo. & Shu mas’ul va qayd ishlatiladi. Har toifa yo‘li uchun mustaqil kutilgan javob borligi isbotlanmagan.\\
\href{https://lex.uz/docs/-6200851}{516-son qaror}, asosiy 6$^{1}$-band, 1-ilova 13 va 44-band & Joriy etishdan oldin zarur ijobiy texnik/himoya xulosasi; birgalikdagi ulanish sinovi; so‘rov soni, davri va statistikasiga kirish. & Qoida va javob sinovi tegishli texnik talab va qabulga qo‘shiladi. Ulanish muvaffaqiyati toifa arifmetikasini isbotlamaydi; tashqi xulosa almashtirilmaydi.\\
\href{https://lex.uz/docs/-8076376}{2026-yil 5-martdagi 93-son qaror}, asosiy 1-band va ilova 3-band & Raqamli texnologiyalar vazirligining infratuzilma va o‘zgarish boshqarishini qamragan IT-audit vazifasi. & Tegishli audit dalili qayta ishlatiladi, takroriy inspeksiya yaratilmaydi. Hujjat Agentlikning muayyan skoring auditi natijasini bermaydi.\\
\href{https://oldmy.gov.uz/oz/service/587}{Rasmiy portal, 587-xizmat} & Mavjud ariza/natija yo‘li, hujjat talab qilinmasligi, bepul xizmat, bir oy deb berilgan muddat va holat sonlari. & Mavjud kanal ishlatiladi. Ko‘rilgan sahifada jami 640\,876, jarayonda 16\,564; yig‘ilish davri, kutayotgan qaror yoshi va ishlov vaqti berilmagan. Bu yillik oqim, kechikkan ish soni yoki bo‘sh quvvat emas.\\
\href{https://gov.uz/oz/news/view/194551}{2026-yil 17-iyuldagi Hukumat xabari} & 184 mahallada bir oydagi 2\,000 asossiz rad va reyestr bo‘yicha Prezidentga murojaat 30\,000 dan 46\,000 ga oshgani bildirilgan; prokuror nazorati topshirilgan. & Ish asosiga kirish va sabab bo‘yicha qayta ko‘rishni qo‘llaydi. Dasturiy sabab, hozirgi uchrash tezligi va taklif ta’sirini isbotlamaydi.\\
\href{https://gov.uz/oz/urganchshahar/news/view/134892}{2026-yil 19-fevraldagi Urganch hokimligi xabari} & Qonuniy kiritish/chiqarish va ma’lumot ishonchliligi bo‘yicha amaliy seminar o‘tkazilgan. & Mos mavjud o‘qitish yo‘li ishlatiladi; mahalliy tadbir butun respublika mutaxassis tayyorligi yoki bo‘sh xodim soatini isbotlamaydi.\\
\bottomrule\end{longtable}\endgroup

\textbf{Qo‘shimcha hissa.} Qo‘shimcha hissa — mustaqil asoslangan javob, sana va chiqarilish natijasiga bog‘langan bajarilgan xulosa va nazorat misollari. Bajarilgan teng sinov qayta ishlatiladi. Bu takror texnik ishni olib tashlaydi, lekin normativ qamrovni o‘zi isbotlamaydi: huquqiy manba alohida tekshiriladi. Ikkala qamrov tasdiqlansa qo‘shimcha norma kerak emas; huquqiy shart yo‘qligi isbotlansa asosiy hujjat yoki qonuniy mavjud vositada vakolatli normativ belgilash kerak. Umumiy vazifa va sun’iy misol har bandning qo‘shimcha foydasini isbotlamaydi.

\textbf{Mutanosib oshkor etish.} 7-band amaldagi uslubiyot va sanani mavjud rasmiy resursda bog‘laydi; 19-band arizachining mavjud yo‘lida asos va e’tirozni beradi; 23-band bayonnomani ichkarida saqlab, mavjud qonuniy hisobotni qo‘llaydi. Chiqarilishdan oldin yoki davriy alohida ochiq hisobot endi talab qilinmaydi. Nashr allaqachon majburiy yoki qonuniy ruxsatli bo‘lsa, ko‘chirma ichki bayonnomadan foydalanib, oshkor qilishga tekshiriladi. Yangi vazifa nashrning o‘ylab topilgan o‘lchangan foydasi bilan asoslanmaydi.

\textbf{Resurs va ijro.} 1-dekabrga kiritishdan oldin ishlab chiquvchi mavjud loyihalash vazifasida huquqiy qaror jadvalini kelishtirib, mavjud qabul hujjatini berilgan misollarga solishtirishi mumkin. 1-yanvardan birinchi huquqiy oqibatli qo‘llanishdan oldin amaldagi norma, zarur texnik/himoya xulosasi va ta’sirlangan yo‘llar sinovi qayd etiladi. Bu bog‘liqlik ketma-ketligi, va’da qilingan muddat yoki boshlanishni kechiktirish vakolati emas. “Inson” xodimiga qo‘shimcha qo‘lda zaxira yuklanmaydi; oddiy ariza qonuniy jarayonini saqlaydi. Sinov joriy etishdan oldingi ruxsatli muhitda, oilaning qo‘shimcha tashrifi yoki hujjatisiz bajariladi.

Haqiqiy qo‘shimcha vazifa uchun ishlab chiquvchi mavjud shartnoma yoki bo‘linma, natija, qo‘shimcha soat, qonuniy mablag‘ va tugash bog‘liqligini belgilaydi. Muddat bo‘yicha amaliy dalil yig‘ish davri, kanal, kirgan/tugagan ish, kutayotgan ish yoshi, ishlov vaqti va manba kechikishini ko‘rsatadi; yig‘ma sayt hisoblagichi buni bermaydi. Har arizaga qo‘shimcha ishlov faqat amaliy dalil va zarur yuklama sinovi bilan joriy etiladi, bo‘sh shtat bor degan bayon bilan emas. Muallifning kichik dasturi tezligi YAMIH o‘tkazuvchanligi deb ko‘rsatilmaydi. 687-son qaror shartli qiymat qoidasini, PQ-388 budjet asosini beradi; loyiha mablag‘ini isbotlamaydi. Taklif nol qo‘shimcha xarajat demaydi va mablag‘siz respublika zaxirasini majburiy va’da qilmaydi. Integratsiyaning aniq qiymati va Agentlikning umumiy muddat imkoniyati o‘lchanmagan; so‘ralgan qaror — topshirilgan loyihadagi tor qayta ishlatish va bo‘shliq vazifasi, haqiqiy kengayish majburiyatdan oldin alohida baholanadi.

\section*{Ilmiy-texnik qism: hisoblash topshirig‘i}
Ushbu qism hisoblash talablarini tushuntiruvchi ishchi texnik materialdir. 2027-yilgi ko‘p o‘lchovli balli modelning aniq koeffitsiyentlari va chegaralari hamda haqiqiy oilalar ma’lumotlari hali olinmagan. Amaldagi 35-son qarorning ochiq hisoblash qoidalari bo‘yicha bajarilgan audit quyida alohida keltirilgan. Umumiy formulalardagi belgilar rasmiy mezon yoki toifa nomlari emas.

Tasdiqlangan qoida umumiy holda \texttt{C = g(F(x), q, v)} shaklida tavsiflanadi: \texttt{x} — tegishli kirish qiymatlari; \texttt{q} — sifat va istisno holatlari; \texttt{v} — qo‘llangan talqin. \texttt{F} chiziqli bo‘lishi shart emas. Vazirlik va idoralarning amaldagi rad etish va yordam tayinlash qoidalari alohida hisobga olinadi.

Faqat qoida qo‘shiluvchi ko‘rinishda bo‘lsa:

\texttt{S = \ensuremath{\Sigma} w\_j z\_j}.

\texttt{z\_j \ensuremath{\in} [a\_j, b\_j]} deb asoslangan yopiq oraliqlar mavjud bo‘lganda:

\texttt{L = \ensuremath{\Sigma} min(w\_j a\_j, w\_j b\_j)};

\texttt{U = \ensuremath{\Sigma} max(w\_j a\_j, w\_j b\_j)}.

Kirish qiymatlari bir-biridan mustaqil ravishda o‘z oralig‘idagi istalgan ratsional qiymatni olishi mumkin bo‘lganda bu oraliq aniq. Diskret yoki o‘zaro bog‘langan kirish qiymatlarida u konservativ tashqi chegara bo‘lib, amalda erishib bo‘lmaydigan toifalarni ham qamrab olishi mumkin. Chegaralar uslubiy yoki ma’lumotlar asosiga ega bo‘lishi kerak. Oraliq bir nechta toifani qamrasa, hisoblash noaniqligi qayd etiladi. Ushbu hisoblash hal etilmagan ma’lumot o‘rniga qiymat tayinlash yoki yordamni rad etish vakolatini bermaydi.

Ishchi dasturda sinov uchun \texttt{C(S) = |\{t: t \ensuremath{\leq} S\}|} ishlatiladi, ya’ni tenglik yuqori indeksli toifaga tegishli. Bu faqat sinov konvensiyasi; rasmiy qoidadagi tenglik, yaxlitlash, istisnolar va yo‘nalish alohida tasdiqlanadi.

Takror hisoblash nomuvofiqligi: \texttt{R = nomuvofiq toifalar soni / taqqoslash mumkin bo‘lgan holatlar soni}.

Sezgirlik ssenariysi uchun: \texttt{F\_s = toifasi almashgan holatlar soni / ssenariyda ikkala toifasi aniqlangan holatlar soni}. Hisobotda ssenariy, surat, maxraj va chiqarilgan yoki hal etilmagan holatlar alohida beriladi. Maxraj nol bo‘lsa, ko‘rsatkich aniqlanmagan deb qayd etiladi.

3-materialga elektron hisoblash qo‘shimchasi additiv namuna, aniq ratsional arifmetika, yopiq intervallar, chegaradagi tenglik va nazorat toifasi bilan solishtirishni beradi. Alohida ochiq qoidalar auditi va mustaqil arifmetik tekshiruv mavjud. Normani baholash uchun zarur xulosalar ushbu tushuntirishda takrorlangan; repozitoriyga kirish shart emas.

\section*{Bajarilgan hisoblash auditi va undan kelib chiqadigan aniq takliflar}
\textbf{Audit sanasi:} 2026-yil 30-sentabr. Aniq kirish qiymatlari, nazorat hisoblari va bajarish natijalari 3-materialning elektron hisoblash qo‘shimchasida beriladi. Asosiy xulosalar quyidagi jadvalda takrorlangan. Bu ochiq qoidalar va tuzilgan holatlar auditi, haqiqiy oilalar kuzatuvi yoki 2027-yil rasmiy modelining tekshiruvi emas.

\begingroup
\fontsize{11bp}{13.2bp}\selectfont
\begin{longtable}{p{\dimexpr\linewidth/3-2\tabcolsep\relax} p{\dimexpr\linewidth/3-2\tabcolsep\relax} p{\dimexpr\linewidth/3-2\tabcolsep\relax}}
\toprule
Tekshiruv & Hisoblangan natija & Normativ oqibat va cheklov \\
\midrule
\endhead
35-son qaror, 1-ilova 30-banddagi rasmiy misol & \texttt{6 480 000 / 5 / 3 = 432 000} so‘m; e’lon qilingan qiymatdan farq \texttt{0} & Misol arifmetikasi takrorlandi. Misoldagi 2025-yil sentabr murojaati uchun may–iyul davri saqlangan; davr umumiy qoidasining kalendar talqini alohida aniqlashtiriladi. \\
32, 33(a), 39(b)-bandlar: daromad aynan \texttt{P} & \texttt{P = 715 000}: \texttt{daromad > P} noto‘g‘ri; \texttt{daromad < P} ham noto‘g‘ri & Boshqa rad etish asosi bo‘lmagan, mehnatga layoqatli yangi arizachi, avvalgi asosiy toifadan chiqarilmagan va og‘ir vaziyat tavsiyasiga ega bo‘lmagan holatda 32-band tan olishni nazarda tutadi, 39-band esa mos toifani bermaydi. Operatsion qarorlar tekshirilmagan. PF-258 3-band ham qat’iy \texttt{< P} shartini qo‘llaydi. \\
27-band yer formulasi: tuzilgan \texttt{(UM, QM) = (100, 90)} m² & \texttt{100 − 90 − min(20, 200) = −10} m² & \texttt{0 \ensuremath{\leq} QM \ensuremath{\leq} UM} algebraik shart bajariladi. Haqiqiy kadastrda bunday holat mavjudligi yoki qo‘shimcha cheklov bilan chiqarib tashlanishi hali aniqlanmagan. \\
Shu yer holatining bir oylik normativ daromadi & \texttt{−10 / 100 \ensuremath{\times} 0,035 \ensuremath{\times} 1 360 000 \ensuremath{\times} 1 = −4 760} so‘m & \texttt{M = 1 360 000} — 2026-yil sentabrdagi rasmiy qiymat; mahalla koeffitsiyenti \texttt{1} — faqat ochiq belgilangan sinov farazi. Taklif etilayotgan nol chegarasi qo‘llansa qiymat \texttt{0}; bu amaldagi formula natijasi bilan aralashtirilmaydi. \\
23-banddagi erkak/ayol koeffitsiyentlari & To‘rt hudud turining har birida farq \texttt{0,4M = 544 000} so‘m/oy; tuzilgan 4 kishilik oilada bir kishiga oylik hissasi \texttt{136 000} so‘m & Bu normativ daromad hissalarining farqi; haqiqiy daromad, kamsitish yoki noto‘g‘ri yordam tayinlashning o‘lchovi emas. Oldingi oylar uchun sentabrdagi \texttt{M} qiymati orqaga tatbiq etilmaydi. \\
Yer formulasining kengaytirilgan nazorat sohasi & \texttt{UM = 100, 1000, 1001, 2000} m² uchun har bir butun \texttt{QM = 0,…,UM}: jami \texttt{4105} holat, shundan \texttt{620} holatda manfiy maydon & Bu tuzilgan holatlar soni, haqiqiy oilalar ulushi emas. \texttt{UM \ensuremath{\leq} 1000} bo‘lsa manfiylik sharti \texttt{QM > 0,8UM}; \texttt{UM > 1000} bo‘lsa \texttt{QM > UM − 200}. \\
Og‘ir vaziyat yo‘lining daromad oralig‘i & 34–35-bandlarning daromad shartlari \texttt{P < D \ensuremath{\leq} 2P} ga teng & Tavsiya, kvota, tegishli oila holati, qayta baholash va boshqa qonuniy talablar saqlanadi. Oddiy \texttt{D} yoki kamaytirilgan \texttt{0,75D} ni umumiy \texttt{[P; 1,5P]} oralig‘iga qo‘yish bu alohida yo‘lni to‘liq takrorlamaydi. \\
24-band soliq asosidagi chegara va 28-band echki hisob-kitobi & Soliqning uch oylik summasi \texttt{T} va oylik baza \texttt{B} uchun \texttt{max(2,5T/3; B)}, kesishish nuqtasi \texttt{T = 6B/5}; 14 echki — \texttt{0}, 15 echki — \texttt{34 000} so‘m/oy & Soliq summasi va oylik baza bir xil davrga keltiriladi. Echki hisobida \texttt{0,07} koeffitsiyent va dastlabki bir shartli bosh istisnosi qo‘llangan; boshqa hayvonlar va operatsion yaxlitlash to‘liq amalga oshirilmagan. \\
\bottomrule
\end{longtable}
\endgroup
\textbf{1. Toifalar chegaralarini muvofiqlashtirish.} Ilovaning 5 va 12-bandlari yuqoridagi tenglik holatini majburiy sinovga kiritadi. Yakuniy siyosiy-huquqiy yechim yuqori yuridik kuchga ega hujjat bilan birga belgilanadi. Texnik sozlama orqali \texttt{<} ni \texttt{\ensuremath{\leq}} ga almashtirish taklif etilmaydi.

Muhokama uchun, tenglikni asosiy toifaga kiritish varianti tanlansa, PF-258-son Farmon 3-bandining (a) va (b) kichik bandlari hamda 35-son qaror 1-ilovasi 39-bandining (a) va (b) kichik bandlaridagi daromadga doir \textbf{“minimal iste’mol xarajatlaridan kam bo‘lgan”} sharti \textbf{“minimal iste’mol xarajatlaridan oshmagan”} deb muvofiqlashtiriladi. PF-258-son Farmon 1-ilovasi izohining (a) kichik bandidagi \textbf{“minimal iste’mol xarajatlaridan kam bo‘lganda”} sharti ham \textbf{“minimal iste’mol xarajatlaridan oshmaganda”} deb o‘zgartiriladi.

PF-258-son Farmon 3-bandining (v) va 35-son qaror 1-ilovasi 39-bandining (v) kichik bandlarida ikki alohida daromad yo‘li belgilanadi. Muhokama uchun taklif etilayotgan mezon matni:

\begin{quote}

“«Kambag‘allik chegarasidagi oila» toifasiga:
«davlat ta’minotidagi oila» yoki «kambag‘al oila» toifasidan chiqarilgan, oilaning har bir a’zosiga to‘g‘ri keladigan bir oylik jami o‘rtacha daromadi, og‘ir hayotiy vaziyat uchun kamaytiruvchi koeffitsiyent qo‘llanmasdan hisoblanganda, minimal iste’mol xarajatlari miqdoridan yuqori, biroq uning bir yarim baravaridan oshmagan oilalar;
og‘ir hayotiy vaziyatga tushgan oilalarni kiritishning alohida tartibiga muvofiq, «mahalla yettiligi»ning tavsiyasi asosida kiritiladigan oilalar ajratiladi. Ushbu yo‘lda 25 foiz kamaytiruvchi koeffitsiyent qo‘llangan daromadning minimal iste’mol xarajatlari miqdorining bir yarim baravaridan oshmasligi tekshiriladi; kamaytirilgan daromad minimal iste’mol xarajatlari miqdoriga teng yoki undan kam bo‘lishining o‘zi kiritishni rad etish asosi bo‘lmaydi. Tegishli oila toifasi, dastlabki baholash, tavsiya, belgilangan miqdor va qayta baholashga doir qonunchilikdagi talablar saqlanadi.”

\end{quote}
Ikkinchi yo‘l PF-258-son Farmon 1-ilovasi izohining (b) kichik bandi va 35-son qaror 1-ilovasining 34–35-bandlari bilan muvofiqlashtiriladi. Oddiy daromad va 25 foiz kamaytirilgan daromad bir xil qiymat sifatida taqqoslanmaydi. Bu mualliflik varianti yordamga huquq va toifaga bog‘liq oqibatlar bo‘yicha ishlab chiquvchining bahosi, Prezident hujjatiga o‘zgartirish va tegishli boshqa normalarni muvofiqlashtirishni talab etadi; faqat Vazirlar Mahkamasi qaroriga o‘zgartirish kiritish bilan hal qilinmaydi. Agar tenglik uchun boshqa qonuniy yo‘l tanlansa, uning toifa, rad etish va istisno shartlari ham ochiq, o‘zaro muvofiq ko‘rsatilishi shart.

\textbf{Taqrizdan keyingi qo‘shimcha nazorat holati (dastlabki holatlar ro‘yxatiga kirmagan):} 3 kishilik, qonuniy tavsiya va mavjud kvota shartlariga mos, boshqa rad etish asosi bo‘lmagan, bolasini yolg‘iz tarbiyalayotgan tuzilgan oilaning uch oylik daromadi \texttt{8 580 000} so‘m bo‘lsin. Oddiy daromad \texttt{8 580 000 / 3 / 3 = 2 860 000 / 3 > 715 000}; kamaytirilgan daromad esa \texttt{(2 860 000 / 3) \ensuremath{\times} 3/4 = 715 000}. 35-bandning (b) kichik bandi bo‘yicha bu kamaytirilgan qiymat yuqori chegaradan oshmaydi. Barcha yo‘llar uchun quyi chegarani \texttt{>P} qilish varianti ushbu holatga zid bo‘lgani sababli qo‘llanilmaydi. Yuqoridagi ikki yo‘lli matn bu muammoni hisobga oladi. Ushbu nazorat huquqiy shartlarni tekshiradi, kuzatilgan oilaning natijasini tasvirlamaydi.

\textbf{2. Yer hisobida manfiy maydonni oldini olish.} 35-son qaror 1-ilovasi 27-bandidagi maydon formulasini quyidagicha bayon etish taklif qilinadi:

\begin{quote}

“TEM = max \{0; UM − QM − min (UM \ensuremath{\times} 20 foiz; 200 m²)\}.
Bunda \texttt{min} qavs ichidagi qiymatlarning eng kichigini, \texttt{max} esa eng kattasini bildiradi. TEM qiymati manfiy bo‘lishi mumkin emas. UM va QM bir xil o‘lchov birligida olinadi. UM yoki QM haqidagi ma’lumot mavjud bo‘lmaganda yoki QM qiymati UM qiymatidan katta bo‘lganda, ushbu ma’lumot vakolatli manbada aniqlashtiriladi; noma’lum yoki o‘zaro zid ma’lumot nol qiymat bilan almashtirilmaydi. Aniqlashtirish ushbu holatning o‘zi asosida yangi rad etish asosini joriy etmaydi.”

\end{quote}
Bu normativ o‘zgartirish taklifi; mavjud dasturda shunday cheklov allaqachon borligi tasdiqlanmagan. Yer solig‘i koeffitsiyentini olish va normativ daromadning boshqa shartlari saqlanadi. Qo‘llash sanasi, qayta hisoblash va oldingi qarorlarga ta’sir tartibi yakuniy hujjatda belgilanadi.

\textbf{3. Davr va versiyani aniqlashtirish.} Ilovaning 4, 6 va 8-bandlari har bir qiymatning oylik yoki davr bo‘yicha summa ekanini, tegishli kalendar oylarini va o‘sha oyda amal qilgan parametrni qayd etishni talab qiladi. 35-son qaror 26-bandidagi pul o‘tkazmasini \texttt{2M} bilan solishtirish davri aniqlanmaguncha migratsiya uchun to‘liq hisoblash qoidasi dasturga kiritilmagan. 19-banddagi valyuta kursining YAMIH ma’lumot olgan kundagi qiymatini qo‘llash qoidasi takror hisoblashda shu sana va kursni saqlash zaruratini tug‘diradi.

Kengaytirilgan audit ikki \textbf{shartli} talqinni alohida hisoblaydi; ulardan birortasi Agentlikning amaldagi algoritmi deb ko‘rsatilmaydi. Solishtirilayotgan \texttt{q} oylik miqdor bo‘lsa, ro‘yxatdagi davlatlar uchun \texttt{q = 2M − 1} holatidan \texttt{q = 2M} holatiga o‘tishda oylik normativ hissa \texttt{3M} dan \texttt{2M} ga, ya’ni \texttt{1 360 000} so‘mga kamayadi. \texttt{q} uch oylik summa deb talqin qilinsa, shu shartli o‘zgarish \texttt{3M} dan \(2M/3\) ga, ya’ni \texttt{9 520 000/3} so‘mga teng. Boshqa davlatlar uchun tegishli shartli kamayishlar \texttt{0} va \texttt{5 440 000/3} so‘mdir. Bu yordam summasi emas. Davr va chegara oldidagi normativ qiymatdan haqiqiy qiymatga o‘tish qoidasi ishlab chiquvchi tomonidan aniq belgilanib, huquqiy va ilmiy jihatdan asoslantiriladi; dasturchi o‘z tashabbusi bilan qo‘shimcha eng kam daromad chegarasini kiritmaydi.

Yer formulasiga nol chegarasi kiritish manfiy normativ hissani oshirishi mumkin. Sentabrdagi \texttt{M} va sinovdagi \texttt{k = 1} uchun oylik uy xo‘jaligi daromadi hissasining eng katta o‘zgarishi deklaratsiya qilingan algebraik sohada \texttt{95 200} so‘mdir; bir kishilik hissa uchun bu qiymat oila a’zolari soniga bo‘linadi. Taklif avtomatik ravishda yordamni ko‘paytiradi deb talqin qilinmaydi. Amaldagi toifa, to‘lov va avvalgi qarorlarga ta’sir qonuniy ma’lumotlar asosida alohida qayta hisoblash orqali baholanadi.

Ushbu natijalar koeffitsiyentlarni bekor qilish, yordam miqdorini qisqartirish yoki ma’lum foizdagi oilalarning toifasi o‘zgaradi degan xulosaga asos bo‘lmaydi. Bunday natijalar uchun rasmiy model, vakolatli muhitdagi qonuniy ma’lumotlar va oldindan belgilangan empirik tekshiruv zarur.

\section*{Amaldagi va taklif etilayotgan yechimlarning qiyosiy bayoni}
Quyida amaldagi qoidaning predmeti qisqa bayon qilinadi; bu qonun matnining so‘zma-so‘z to‘liq ko‘chirmasi emas. Takliflar yakuniy hujjatning darajasi va rasmiy tahriri bilan muvofiqlashtiriladi.

\textbf{1. Toifa va tenglik.} Amaldagi manbalar: PF-258 3-band va 1-ilova izohlari; 35-son qaror 1-ilovasi 32–35 va 39-bandlar. Taklif: yuqoridagi ikki yo‘lli toifa matni yoki ishlab chiquvchi tanlagan boshqa o‘zaro izchil tenglik yechimi. Asos: daromad aynan \(P\) bo‘lganda tan olish va toifaga ajratish shartlarining bir xil natija bermasligi; og‘ir vaziyat yo‘lining alohida huquqiy shartlari saqlanishi. Bu variant asosiy tekshiruv ilovasini qabul qilishning majburiy sharti emas.

\textbf{2. Yer formulasi.} Amaldagi manba: 35-son qaror 1-ilovasi 27-band. Taklif: yuqoridagi nol chegarasi va bir xil kadastr obyektiga tegishli maydonlar uchun tekshiriladigan manba qoidasi. Asos: deklaratsiya qilingan algebraik sohada manfiy maydon. Nol chegarasi qo‘llanishidan oldin hisoblangan daromad oshishi mumkin bo‘lgan holatlar va oldingi to‘lovlarga ta’sir baholanadi; geometrik ma’noga mos boshqa yechim tanlansa, u qayta hisoblanadigan tarzda bayon qilinadi.

\textbf{3. Natija va hisoblash asoslari.} Amaldagi manba: 35-son qaror 1-ilovasi 37-banddagi SMS va tegishli qaror qoidalari. Taklif etilayotgan qo‘shimcha norma: “Qaror natijasi bilan bir vaqtda ariza beruvchiga qo‘llangan hisoblash talqini, asosiy qiymatlar va davrlar, normativ qiymat yoki istisno sababi, qarorning huquqiy va faktik asoslari hamda e’tiroz bildirish yo‘li elektron kabinet orqali taqdim etiladi. Elektron kabinetdan foydalanmaydigan shaxs mazkur axborotni «Inson» markazi orqali olishi mumkin. SMS-xabar mazkur axborotning o‘rnini bosmaydi.” Asos: yangi ilovaning 18–20-bandlari; shaxsga doir ma’lumotlar haqidagi Qonun 24-modda va tegishli ma’muriy tartiblar. Himoyalangan uchinchi shaxs ma’lumoti chiqarilmaydi.

\textbf{4. Hisoblashga e’tiroz va kengash vakolati.} Amaldagi manba: 35-son qaror 4-ilovasi 26 va 31-bandlar kengash predmetini “mahalla yettiligi” hujjatlari bilan chegaralaydi; 33-band shikoyatni o‘n besh kunda ko‘rishni belgilaydi. Avtomatik hisoblash e’tirozi uchun kengashga yangi umumiy vakolat bor deb talqin qilinmaydi. YAMIH va “Inson” markazidagi qabul shakliga taklif: “Faqat avtomatlashtirilgan ishlov berishga asoslangan qaror yuzasidan e’tiroz ma’lumotlar mulkdori yoki operatoriga yuboriladi va uning natijasi haqida subyektga o‘n kun ichida yozma xabar beriladi. «Mahalla yettiligi» hujjatiga shikoyat apellyatsiya kengashiga tegishli predmeti bo‘yicha yuboriladi. Bir murojaatdagi turli masalalar bo‘yicha dastlabki kelib tushish sanasi va tegishli yo‘naltirishlar saqlanadi.” Asos: O‘RQ-547 24-modda. Manba ma’lumotini subyektning so‘roviga binoan tuzatish yoki to‘ldirish 11-modda bo‘yicha uch kun ichida, aniqlangan noto‘g‘ri ma’lumotni tuzatish esa darhol amalga oshiriladi; bu muddat o‘n kunlik avtomatik qarorga e’tiroz muddatidan alohida hisoblanadi. Kengash vakolatini kengaytirish tanlansa, uning predmet, dalil talab qilish, qaror va ijro normalari alohida muvofiqlashtirilishi kerak; bu taklif bunday kengaytirish allaqachon amalga oshganini da’vo qilmaydi.

\textbf{5. Tizim uzilishi.} Amaldagi manba: 35-son qaror 1-ilovasi 48-banddagi vaqtinchalik algoritmni o‘zgartirish vakolati. Qo‘shimcha norma taklifi: “Vaqtinchalik talqinning huquqiy asosi, doirasi, boshlanishi, tiklash sharti va qo‘llangan hisob-kitoblar qayd etiladi. Tizim tiklanganda vaqtinchalik talqin qo‘llangan holatlar aniqlanadi va tegishli qonuniy qoidalar bo‘yicha tekshiriladi.” Asos: uzilish vakolatini yashirin doimiy mezon almashtirishga aylantirmaslik; yangi ilovaning 11, 23 va 26-bandlari. Ushbu qo‘shimcha mavjud vakolatni kengaytirmaydi.

\textbf{6. Axborot almashish va to‘lovlar.} Muvofiqlashtirish predmeti: 35-son qarorning tegishli axborot almashish ilovalari va 3-ilovadagi tayinlash hamda to‘lov qoidalari. Taklif: manba davri, olingan vaqt, sifat holati va qaror talqinini uzatish maydonlariga qo‘shish; qonuniy tayinlangan yordamning avtomatik to‘xtashini yangi “hal etilmagan” texnik holat bilan bog‘lamaslik. Ilgari hisoblangan daromadni tuzatish, yangi siyosiy mezon va to‘lovni undirish bir xil hodisa deb belgilanmaydi. Aniq to‘lov bandlari rasmiy 2027-yilgi loyihaning tanlangan mezonlari bilan solishtirilgach ishlab chiquvchi tomonidan rasmiy o‘zgartirish jadvaliga kiritiladi; hozir bu jadval barcha yordam dasturlarini qamrab olgan deb ko‘rsatilmaydi.

\textbf{7. Amaldagi mas’uliyat va qo‘shimcha tekshiruv tartibi.} 35-son qarorning 1-ilovasi 46-bandida Reyestr va YAMIHning to‘g‘ri ishlashi, modulni takomillashtirish va axborot almashish uchun Agentlik huzuridagi “Axborot texnologiyalar markazi” davlat muassasasi mas’ullarining, o‘z vaqtida va ishonchli javoblar uchun tegishli vazirlik hamda idora rahbarlarining javobgarligi belgilangan. 49-band nazoratni Bosh prokuratura, Agentlikning ichki audit xizmati va Ijtimoiy inspeksiyaga yuklaydi. Taklif yangi inspeksiya tuzmaydi va ularning qonuniy vakolatini boshqa organga o‘tkazmaydi. U shu mas’uliyat uchun tekshirish dalillarini aniqlashtiradi: mustaqil takrorlanadigan nazorat hisobi, har bir talqinni qabul qilish, davr va istisnolarni saqlash, hal etilmagan holatlar soni, qaror asoslari va oshkor qilish mumkin bo‘lgan umumlashtirilgan natijalar. Mazkur ikki bandning o‘zida shu batafsil bayonnoma belgilanmagan. Bu aniqlangan normalar bilan qiyoslash, butun qonunchilik yoki amaldagi tizimda teng mazmundagi vazifa va vositalar yo‘qligining dalili emas. Ishlab chiquvchi mavjud teng yechimlarni aniqlashi, ularning vakolatini saqlashi va takrorlanadigan qoidalarni birlashtirishi lozim. Teng yechimning mavjudligi takrorlanishdan qochish uchun asos bo‘la oladi, alohida talabni javobsiz qoldirish uchun emas.

\subsection*{Asosiy tekshiruv ilovasi: amaldagi vazifalar bilan bog‘lanish}
Ushbu qiyos 1-materialning majburiy qismiga taalluqli. Birinchi ustunda aniq amaldagi norma predmeti, ikkinchisida muallif taklif qilgan qo‘shimcha ijro vositasi berilgan. Bu butun qarorning so‘zma-so‘z taqqoslash tahriri emas, tahliliy bog‘lanishdir. Ikkinchi ustundagi bandlar 1-material ilovasiga tegishli.
\begingroup\fontsize{11bp}{13.2bp}\selectfont
\begin{longtable}{p{\dimexpr\linewidth/3-2\tabcolsep\relax} p{\dimexpr\linewidth/3-2\tabcolsep\relax} p{\dimexpr\linewidth/3-2\tabcolsep\relax}}
\toprule
Amaldagi vazifa & Taklif etilgan ijro & Tekshiriladigan zarurat\\
\midrule\endhead
35-son qaror 1-ilovasi 46-bandi: Reyestr va YAMIHning qonunchilik va texnik hujjatlarga muvofiq to‘g‘ri ishlashi. & 4–7, 12–17, 25–26-bandlar: bog‘langan qoidalar tavsifi, mustaqil kutilgan javoblar, qismlar tekshiruvi va aniq talqinni qabul qilish. & Belgilangan misollarda tasdiqlangan qoidaga muvofiqlikni aniqlash; mos jami yoki ikki dastur qismlar xatosi yoxud umumiy noto‘g‘ri javobni yashirmaydi.\\
\midrule
45–46-bandlar: ma’lumotlarni himoya qilish va davlat manbalarining o‘z vaqtida ishonchli javoblari. & 8–10, 21–22, 28-bandlar: alohida sifat holatlari, vakolatli aniqlashtirish, qonuniy ruxsat, berilgan kod tekshiruvi va saqlash chegarasi. & Kelmagan javob nolga aylanmaydi; davlat manbasi vazifasi oilaga o‘tkazilmaydi va himoyalangan kirishlar tashqi tekshiruv xizmatiga berilmaydi.\\
\midrule
37-band: natija haqida SMS. Shaxsga doir ma’lumotlar qonuni 24-moddasi: avtomatik qarorni tushuntirish va e’tiroz natijasini o‘n kunda xabar qilish. & 18–20-bandlar: hisob asoslariga kirish, manba, dastur va qoidani ajratish, dastlabki qabul qaydi va vakolatli yo‘naltirish. & Oila tegishli sababni aniqlab nizolashi mumkin; manba tuzatish mahalla hujjati ustidan shikoyat bilan almashtirilmaydi.\\
\midrule
48-band: almashish uzilgandagi vaqtinchalik vakolat. 49-band: sanalgan organlarning o‘z vakolatidagi nazorati. & 11, 17, 23–24, 27-bandlar: vaqtinchalik qayd, ta’sir doirasi, tiklash, ichki bayonnoma va avvalgi qarorlarni qonuniy qayta ko‘rish. & Vaqtinchalik hisoblar aniqlanadigan bo‘lib qoladi; texnik tuzatish to‘lovni kamaytirish yoki undirish uchun mustaqil asos yaratmaydi.\\
\bottomrule\end{longtable}\endgroup
Tenglik va yer formulasi haqidagi ixtiyoriy variantlarni chiqarish ushbu tekshiruv vazifalarini bekor qilmaydi. Teng amaldagi tartib dalil bilan aniqlansa saqlanadi va takrorlash birlashtirilganda uning normasi ko‘rsatiladi. Ushbu qiyos 2027-yil ehtiyoj mezonlarini tasdiqlamaydi.

\subsection*{Tanlangan normalar bo‘yicha uch grafali taqqoslash jadvali}
Quyidagi jadval O‘RQ-682 ilovasi 85-bandining uch grafali tuzilishidan foydalanadi. Birinchi grafada faqat o‘zgartirish tegishli bo‘lgan ibora yoki hisoblash qismi keltirilgan; boshqa shartlar bekor qilinmaydi. Bu mualliflik taklifidagi tanlangan normalar jadvali bo‘lib, rasmiy ishlab chiquvchining barcha hujjatlar va yordam dasturlarini qamrab oladigan yakuniy jadvali emas.
\begingroup\fontsize{11bp}{13.2bp}\selectfont
\begin{longtable}{p{\dimexpr\linewidth/3-2\tabcolsep\relax} p{\dimexpr\linewidth/3-2\tabcolsep\relax} p{\dimexpr\linewidth/3-2\tabcolsep\relax}}
\toprule
Amaldagi tahrirning tegishli qismi & Taklif etilayotgan tahrir & O‘zgartirishning asosi\\
\midrule\endhead
PF-258 3-band (a), (b), 1-ilova izohi (a) va 35-son qaror 1-ilovasi 39-band (a), (b): daromad uchun “\textit{\uline{minimal isteʼmol xarajatlaridan kam}}” sharti. & Tanlangan inklyuziv variantda aynan shu daromad iborasi “\textbf{minimal iste’mol xarajatlaridan oshmagan}” deb almashtiriladi. Prezident va Vazirlar Mahkamasi darajasidagi barcha ko‘rsatilgan qismlar birgalikda o‘zgartiriladi; boshqa shartlar saqlanadi. & Daromad aynan \(P\) bo‘lgan yangi ariza bo‘yicha tanlangan shartlarda tan olish va asosiy toifa predikatlari mos emas. Bu variant siyosiy tanlov bo‘lib, alohida xarajat va tatbiq etish bahosini talab qiladi.\\
\midrule
35-son qaror 1-ilovasi 35-band (b): maxsus yo‘lda 25 foiz kamaytirilgan daromadning \(1.5P\) dan oshmasligi; tavsiya, kvota va boshqa shartlar mavjud. & \textbf{Oddiy daromad yo‘li bilan maxsus kamaytirilgan daromad yo‘li ajratilsin. Maxsus yo‘lda \(A\leq1.5P\) tekshirilsin; \(A\leq P\) bo‘lishi o‘zi rad etish asosi bo‘lmasin. Boshqa qonuniy shartlar saqlansin.} & Barcha yo‘llar uchun \(A>P\) degan yangi quyi chegara tuzilgan uch kishilik nazorat holatida \(A=P\) ga zid. Maxsus yo‘lni umumiy interval bilan qisqartirish mumkin emas.\\
\midrule
35-son qaror 1-ilovasi 27-banddagi tegishli maydon hisobining algebraik ko‘rinishi: \textit{\uline{\(L=U-B-\min(U/5,200)\)}}. & Ishlab chiquvchi geometrik ma’noni tasdiqlagan holda taklif etilayotgan hisob: \textbf{L = max\{0, U − B − min(U/5, 200)\}}. Maydonlar bir obyekt va davrga tegishli bo‘lsin; nol chegarasi, manba va o‘lchov birligi normativ qoidadan aniqlansin. & \(0\leq B\leq U\) algebraik sohada manfiy natija bor. Nol chegarasi hisoblangan daromadni oshirishi mumkin; u yordamni avtomatik oshirish yoki oldingi to‘lovni undirish asosi emas. Operatsion uchrash chastotasi aniqlanmagan.\\
\midrule
35-son qaror 1-ilovasi 37-band: “Oilaning Reyestrga kiritilganligi yoki Reyestrga kiritish rad etilganligi haqida ariza beruvchining mobil telefon raqamiga SMS-xabar yuboriladi.” & Mazkur SMS talabi saqlansin. \textbf{Qaror bilan birga hisoblash talqini, asosiy qiymatlar va davrlar, huquqiy va faktik asoslar hamda e’tiroz yo‘li elektron kabinetda ochilsin; kabinetdan foydalanmaydigan shaxs axborotni “Inson” markazidan olsin. Himoyalangan uchinchi shaxs ma’lumoti chiqarilmasin.} & Natija xabari hisoblash asosi va e’tiroz tartibini to‘liq ochmaydi. Taklif shaxsga doir ma’lumotlar haqidagi Qonun 24-moddasi va tatbiq etiladigan ma’muriy tartiblar bilan muvofiqlashtiriladi.\\
\bottomrule\end{longtable}\endgroup

\section*{Muhokama, kelishish va havolalarning asosi}
Ushbu mualliflik taklifi davlat organlariga yuborilmagan, davlat portalidagi jamoatchilik muhokamasidan o‘tmagan va manfaatdor idoralar bilan rasmiy kelishilmagan. Ichki mustaqil tekshiruv belgilangan doirada hisoblash sifatini qo‘llab-quvvatlaydi, idoralar kelishuvi emas. Idoralarning roziligi yoki qarshiligi haqida dalil yo‘q; e’tirozlarga javoblar — muallif ko‘rib chiqqan ehtimoliy masalalar.

22-modda asosiy 4-banddagi tayyorlashda ishtirok, ilmiy tavsiya va mutaxassis maslahatlarini asoslaydi. PF-193 topshiriq va mavzuni belgilaydi. PF-258 va 35-son qaror toifalar, hisoblash, murojaat va vaqtinchalik vakolatning amaldagi qoidalarini ko‘rsatadi. Bu havolalar quyi hujjat yoki texnik yo‘riqnomaga yuqori huquqiy kuchdagi mezonni almashtirish vakolatini bermaydi. Shaxsga doir ma’lumotlar va ma’muriy tartiblar haqidagi qonunlar ma’lumot nazorati, qaror asoslariga kirish va qonuniy e’tiroz yo‘lini asoslaydi.

\section*{Ijro ketma-ketligi va javobgarlik}
Uslubiyot vakolatli tasdiqlangach, mavjud uslubiy, tizim va yuridik vazifalar birlik, davr, chegara, istisno, sana va asosni muvofiqlashtiradi. Mustaqil misol amaldagi texnik jarayon va chiqarilishga kiritiladi; teng sinovga havola qilinadi. 25–26-band qabuli majburiy tashqi xulosani almashtirmaydi. Tavsif mavjud resursda bog‘lanadi, bayonnoma 23-bandda ichkarida saqlanadi. Nashr mavjud vazifa yoki ruxsatga ergashadi; yangi ochiq hisobot ishga tushish sharti emas. Qonuniy muddat va yanvar boshlanishi saqlanadi.

Vakolatli manba egasi yoki operatori fakt ma’lumotini, tizim bo‘linmasi dastur nuqsonini hal etadi, uslubiy bo‘linma yuridik xizmat bilan huquqiy tahlil tayyorlaydi. Rasmiy sharh yoki o‘zgartirish kerak bo‘lsa masala vakolatli qabul qiluvchi organga yuboriladi; 55-modda sharh bilan o‘zgartirishni ruxsat etmaydi. Individual qarorni vakolatli subyekt o‘zgartiradi. Texnik qabul siyosat mezonini tasdiqlamaydi. Topshiriq va qonuniy muddat ichki reja bilan ko‘chmaydi.

\section*{Ijro bosqichlari va yakunlanganlik dalili}
\begingroup\fontsize{11bp}{13.2bp}\selectfont
\begin{longtable}{p{0.22\linewidth} p{0.25\linewidth} p{\dimexpr0.53\linewidth-6\tabcolsep\relax}}
\toprule
Vazifa va bandlar & Mas’ul & Material va tugash sharti\\
\midrule\endhead
Uslubiyot va o‘zgarish; 4–7, 12 & Agentlik Vazirlik bilan; uslubiy bo‘linma va yuridik xizmat & Qoida pasporti va o‘zgarish jadvali. Asos, sana, ustuvorlik va barcha tegishli yo‘llar mos; dastur yetishmayotgan qoidani kiritmaydi.\\
\midrule
Manba va sifat; 8–10 & Axborot texnologiyalar markazi mas’ul xodimlari; 46-band doirasida manba idoralar rahbarlari & Manba xaritasi va sifat jurnali. Ishlatilgan maydonning davri, manbasi, holati va aniqlashtirish yo‘li mavjud; nol noma’lumdan farq qiladi.\\
\midrule
Nazorat dasturi; 12–17 & Tizim bo‘linmasi va dastur mas’ulidan boshqa kutilgan javob tekshiruvchisi & Qoida va misol ro‘yxati, aniq javob va bayonnoma. Bajarilgan tarkib e’lon qilinganga mos; ta’sirli farq hal etilib qayta sinalgan.\\
\midrule
Izoh va e’tiroz; 18–20 & Vakolatli “Inson” xodimi, ma’lumot egasi yoki operatori va vakolatli qayta ko‘rish organi & Ochiq izoh namunasi va bog‘langan yuborish qaydi. Oluvchi, muddat va natija ma’lum; yo‘lni tekshirish qonuniy muddatni kechiktirmaydi.\\
\midrule
Uzilish va tiklanish; 11, 23 & Vakolatli mansabdor va tizim mas’ullari & Mavjud vaqtinchalik qayd va tiklangan hisob; qarorlar topiladi, farq qonuniy hal etiladi; ichki dalil yangilanadi.\\
\midrule
Resurs va qabul; asosiy 6–7, 25–26 & Agentlik Vazirlik bilan; rahbar yoki vakolatli o‘rinbosar & Chiqarilish qaydi, tashqi xulosa, sinov, haqiqiy qo‘shimcha vazifa/mablag‘; yangi qo‘lda xizmat quvvati taxmin qilinmaydi.\\
\midrule
Saqlash va qonuniy oshkor; 21–24, 28 & Operator, himoya va mavjud nazorat & Qonuniy jadval, ichki bayonnoma va majburiy/ruxsatli oshkor tekshiruvi; yangi nashr davri yo‘q.\\
\bottomrule\end{longtable}\endgroup

Asosiy loyiha 2026-yil 1-dekabrgacha kiritilishidan oldin qonuniy qaror jadvali va nazoratning mavjud qabul daliliga bog‘lanishi kelishiladi. 2027-yil 1-yanvardan qoida ilk qo‘llanishidan oldin vakolatli tasdiq, amaldagi texnik/himoya xulosasi va yo‘l sinovi qayd etiladi. Bu bog‘liqlik ketma-ketligi, baholangan xodim jadvali yoki boshlanishni kechiktirish vakolati emas. Dastlabki misol loyihada tayyorlanadi; qat’iy javob tasdiqlangan qoidaga ergashadi. Qo‘shimcha integratsiya xarajat majburiyatidan oldin vazifa va mablag‘ dalilini talab qiladi.

\section*{Dastlabki tartibga solish ta’siri va moliyaviy asos}
\textbf{A0 — amaldagi tartib.} Nazorat va qonuniy huquq saqlanadi. Huquqiy mavjudlik aniqlangan; misol qamrovi uchun Agentlik chiqarilish dalili olinmagan. Bu nazorat ishlamasligi xulosasi emas.

\textbf{A1 — huquqiy asosning texnik ijrosi.} Qonuniy qoida alohida aniqlanganda xulosa va misol mavjud talab va chiqarilish qaydiga kiritiladi. Teng sinov va kanal qayta ishlatiladi; qoplanmagan ish va resurs hujjatlanadi. A1 afzal texnik ijrodir, yo‘q normativ asosning o‘rnini bosmaydi.

\textbf{A2 — topshirilgan hujjatning minimal normativ mazmuni.} PF-193 loyihasiga yangi uslubiyotning zarur sharti va aniq mavjud norma ta’minlamagan kafolat kiritiladi. 1-materialning besh qisqa bandi, teng tahrir yoki saqlangan tatbiq norma tanlanadi; A1 texnik ijroni beradi. To‘liq shakl bo‘linmas emas. Alohida qo‘lda xizmat, yangi statistik hisobot va yordam huquqi bo‘yicha Prezident o‘tishi taklif etilmaydi. Qiyos O‘RQ-682 35-modda va PQ-5025 2(b) rasmiy hisoboti emas.

Faqat aniqlangan qo‘shimcha vazifa baholanadi: huquqiy muvofiqlik, yetishmagan misol, tegishli integratsiya/kirish o‘zgarishi va haqiqiy xizmat. Mavjud ish va shartnoma qamrovi takroriy xarid va ikki marta hisobga qarshi qayd etiladi. Agentlik hajm, soat, amaliy dalil va shartnoma doirasini beradi; Vazirlik qonuniy mablag‘ni belgilaydi. Integratsiya, ortiqcha xodim, yakuniy qiymat yoki tejam o‘lchanmagan. Mavjud dastur va kod nol xarajatni isbotlamaydi.

Haqiqiy ehtiyojni empirik tekshirish uchun alohida, qonuniy ma’lumotlar bilan tadqiqot talab etiladi. Ushbu tadqiqotning xarajati va natijasini asosiy hisoblash muvofiqligi sinovidan ajratish mumkin; tugallanmagan empirik tadqiqot yordamni to‘g‘ri nishonlash isboti deb ko‘rsatilmaydi.

\section*{Qabul qilish va mablag‘ uchun zarur dalillar}
Tekshiruv ilovasi ixtiyoriy tenglik va yer formulasi tahrirlaridan alohida qaror qilinadi. U qoida va misol bog‘lanishi ro‘yxati, kutilgan javoblarning mustaqil tekshirilgan huquqiy chiqarilishi hamda tarkibiy qiymatlarni qiyoslashni talab qiladi. Sinov ro‘yxati qo‘llanadigan istisnoni tushirib qoldirib, o‘z to‘liqligini tasdiqlay olmaydi. Yetishmaydigan yoki takroriy holat, hal etilmagan tarkibiy tafovut va nol maxraj ko‘rsatiladi. Bir xil jami bir-birini qoplaydigan xatolarni yashirmaydi. Uslubiyot, dasturiy chiqarilish va sanali sozlama alohida aniqlanadigan qaydlarga ega; o‘sha huquqiy qoida bo‘yicha dastur tuzatishi qayta tiklanadi.

Bajarilgan ish 37 regressiya sinovi va qat’iy misollar qamrovini tekshiruvchi dasturni taqdim etadi. Bu amaliy qabulni bajarmaydi: ilmiy ishning 5-jadvali har qoida guruhi uchun bajarilgan tekshiruv va zarur qo‘shimcha rasmiy manba, yo‘l va yaxlitlash dalilini ko‘rsatadi. O‘nta asosiy daromad qatori va 4\,105 yer qatori — tuzilgan nazoratlar. Ular oila deb hisoblanmaydi va fiskal ta’sir prognozida ishlatilmaydi.

Ishlab chiquvchining tayyorlik bayonnomasi uslubiyot, axborot tizimi va yuridik xizmat, ularning qabul qilingan materiallari, amal sanalari va hal etilmagan masalalarni ko‘rsatishi kerak. Resurs hisobi integratsiya, nazorat, saqlash, kirish, ta’lim va ko‘rib chiqish uchun ish miqdori, qo‘shimcha pul, xodim soati, mavjud quvvat va mablag‘ni qayd etadi. Doimiy va o‘zgaruvchi xarajat ikki marta hisoblanmaydi. Budjet summasi yoki tejam o‘ylab topilmagan. Keyingi farovonlik tadqiqoti alohida baholanadi; quvvat yetishmovchiligi arizachi muddatini uzaytirishga asos emas.

Taklif etilgan nazorat tuzilishining texnik manbalari: \href{https://discovery.ucl.ac.uk/id/eprint/1471263/1/06963470.pdf}{Barr va boshqalar, 2015, sinov javobi haqidagi sharh} va \href{https://nvlpubs.nist.gov/nistpubs/ir/2021/NIST.IR.8397.pdf}{NIST IR 8397, 2.6–2.8-bo‘limlar}. Bular muhandislik manbalari, O‘zbekistonda qo‘shimcha huquqiy majburiyat emas.

\section*{Kutilgan natija, oqibat va tekshiriladigan ko‘rsatkichlar}
Bevosita kutilgan natija — har qo‘llangan talqinning tavsiflangan qoidasi, sinalgan nazorat to‘plami va saqlangan individual qaror asosi. Ijro tegishli mezonlar pasporti, e’lon va bajarilgan sinov tarkibi mosligi, hal etilmagan komponent farqi, izohga kirish va tuzatish hamda e’tirozning qonuniy muddatlari bilan baholanadi. Ulushning surat, maxraj va chiqarilgan holatlari beriladi. Bular tartib ijrosini tekshiradi, ehtiyojni aniqlash aniqligini emas. Ixtiyoriy xato, barqarorlik yoki shikoyat kamayishi foizi belgilanmaydi.

Oila uchun ishlatilgan qiymat va huquqiy asosga kirish, vakolatli manbada tuzatish va faqat texnik holatdan yangi rad asosi yaratilmasligi kutiladi. Organlar bog‘langan nazorat, izoh va hisobotni yuritib, zarur quvvatni ta’minlaydi. Xatoni tuzatish qonuniy natijani har ikki yo‘nalishda o‘zgartirishi mumkin; to‘lovni o‘zgartirish, qayta ko‘rish va undirish o‘z huquqiy asosi va vakolatli individual qarorni talab qiladi. Amaldagi mustaqil rad, Reyestrdan chiqarish va to‘lovni to‘xtatish asoslari saqlanadi. Tuzilgan holat bunday o‘zgarish soni yoki budjet qiymatini ko‘rsatmaydi.

Amaliy xavf — to‘liqsiz manba, xato nazorat javobi, tayyorlanmagan e’tiroz yo‘li, xodim yetishmasligi va kichik guruh orqali shaxsni aniqlash. Mos choralar: alohida holat va aniqlashtirish (8–10), mustaqil huquqiy chiqarish va komponent sinovi (12–17), izoh va yo‘naltirish (18–20), himoyalangan ruxsat va nashr (21–23), resurs va yozma qabul (asosiy 6–7, 25–26). Nazorat dasturining o‘zi xavfni tugatmaydi; tugash dalili ijro jadvalida berilgan.

\section*{Moliyaviy-iqtisodiy hisobning tarkibi}
Asosiy loyiha hisobi qoida tavsifi va huquqiy moslashtirish; manba integratsiyasi va sifat maydoni; mustaqil javob va sinov; talqin qaydi va himoyalangan saqlash; kabinet va “Inson” orqali murojaat; xodim tayyorlash; hisobot va oshkor qilish nazorati; xizmat ko‘rsatish va o‘zgarishni qayta sinashni qamraydi. Har ish uchun bir martalik va yillik miqdor, birlik, birlikka qo‘shimcha pul xarajati, xodim soati, mavjud quvvat, zarur ortish va moliya manbasining hujjati qayd etiladi. Shtat mavjudligi bo‘sh quvvat emas; qayta taqsimlangan mehnat yangi pul xarajatidan ajratiladi.

3-materialdagi bir xil hisob ta’riflari qo‘llanadi:
\[K_{\rm first}=K_{\rm setup}+K_{\rm annual},\qquad K_{\rm annual}=K_{\rm fixed}+\sum_r Q_r k_r.\]
\(Q_r\) — ko‘rsatilgan birlikdagi yillik miqdor, \(k_r\) — birlikka qo‘shimcha pul xarajati; doimiy va o‘zgaruvchan qism takror hisoblanmaydi. Agentlik miqdor va quvvatni o‘z yozuvlaridan, Vazirlik qonuniy mablag‘ manbasi va budjet rasmiylashtirishini belgilaydi. Mustaqil ehtiyoj tadqiqoti va ixtiyoriy chegara tuzatishi alohida hisoblanadi. Yakuniy loyiha qiymati aniqlanmagan: tasdiqlangan ish tarkibi, modullar soni, davomiylik va ajratilgan mablag‘ olinmagan; e’lon qilingan qiymat hisobi qoidalari quyida ko‘rsatilgan. Nol xarajat yoki aniq qo‘shimcha mablag‘ da’vo qilinmaydi; asosiy 6-band uchun zarur dalillar tarkibi beriladi.

\subsection*{Qiymatni aniqlashga oid tekshirilgan idoraviy qoidalar}
\href{https://lex.uz/docs/-6720230}{Vazirlar Mahkamasining 2023-yil 28-dekabrdagi 687-son qarori} amaldagi tahririning asosiy 3-bandi va 3-ilovasi Agentlik Axborot texnologiyalar markazi Ijtimoiy himoya davlat jamg‘armasi yoki Agentlikning budjetdan tashqari mablag‘lari hisobidan bajaradigan loyihalar qiymatini aniqlash tartibini belgilaydi. Tatbiq etilishi ijrochi va mablag‘ manbaiga bog‘liq; bu barcha tekshiruv ishlariga umumiy tarif emas.

Axborot tizimlarini texnik qo‘llab-quvvatlash va takomillashtirish uchun 3-ilovaning 4-bandi mehnatga haq to‘lash fondini belgilaydi:
\[\mathrm{MHTF}=1{,}85\times\mathrm{MS}\times16\times\mathrm{MHTEKM}\times\mathrm{LAOOS}.\]
MS — modullar soni; LAOOS — loyiha amalga oshiriladigan oylar soni; MHTEKM — belgilangan eng kam mehnat haqi, Markaz ijtimoiy soliq to‘lovchi bo‘lsa, shu soliq hisobga olinib oshiriladi. Loyihaga oid boshqa xarajatlar mehnatga haq to‘lash fondining 5 foizi etib belgilangan. 3 va 7–9-bandlar soliq va majburiy to‘lovlar, Markaz, boshqaruv/yordamchi xodimlarni saqlash va foydani alohida ko‘rsatadi; fond formulasi shartnomaning to‘liq narxi emas. 4–6-band loyihalarida 8-band Markaz va boshqaruv/yordamchi xodimlar xarajatini shu bandlar bo‘yicha aniqlangan loyiha qiymatining 40 foizigacha cheklaydi; shtat va haq shartlari Agentlik direktori buyrug‘ida belgilanadi. 9-band Markaz foydasini loyiha xarajatining 2 foizi deb belgilaydi. Foizlar ko‘rsatilgan baza va loyiha turiga qo‘llanadi, barcha ishga farqlamasdan qo‘shilmaydi. Yangi loyiha qiymati 10-band bo‘yicha texnik talab, muddat, mutaxassislar toifasi va sonidan kelib chiqib, Agentlik direktori buyrug‘i asosida aniqlanadi. 12-banddagi shartli shtat va ish haqi parametrlari faqat qiymat hisobi uchun, haqiqiy shtat yoki bo‘sh quvvat dalili emas. Hisobdan oldin ish turi belgilanadi; modullar soni yoki muddat bu yerda taxmin qilinmaydi.

\href{https://lex.uz/docs/-7959569}{2025-yil 26-dekabrdagi PQ-388-son qaror}ning 19–21-bandlari va 6-ilovasi Agentlikda dasturiy budjetlashtirish asosini belgilaydi. Bu budjet sharoiti, mazkur tekshiruv moduliga mablag‘ ajratilgani emas. Agentlik va Vazirlik yakuniy hisobni tegishli tasdiqlangan dastur, smeta va moliyalashtirish qarori bilan bog‘laydi. Ijtimoiy himoyaga umumiy xarajat taklif etilgan hisob tekshiruvi narxi o‘rnida olinmaydi.

\section*{Xalqaro materiallar va to‘plam tillari}
Taklif xorijiy huquqiy normani ko‘chirmaydi va yangi xalqaro shartnomani bajarmaydi. OECD/JRC, Barr nazorat javobi sharhi va NIST tavsiyalari noaniqlik va dastur tekshiruvi uchun muhandislik manbalaridir; huquqiy hisobga moslashtirish 3-materialda asoslanadi. Ular majburiy O‘zbekiston koeffitsiyenti yoki chegarasini bermaydi. Ushbu qonunosti taklifi uchun chet material 84-band doirasida tegishli hajmda ochiladi; xorijiy qonunchilik modeli isbotlangan optimal yechim deb olinmaydi.

Murojaat va 1–3-materialning asosiy tili ruscha. \href{https://lex.uz/docs/-3336169}{Murojaatlar haqidagi qonun} 6-moddasi davlat va boshqa tillarda taqdim etishga ruxsat beradi. O‘zbekcha tahrir \href{https://lex.uz/docs/-5378966}{O‘RQ-682} 21-moddasi bo‘yicha rasmiy loyihani davlat tilida tayyorlashga beriladi; inglizcha mos tarjimadir. Uchala matnda formula, tengsizlik, muddat va vakolat solishtiriladi. Mazmuniy tarjima farqi qulay matn tanlash bilan emas, ishlab chiquvchi yakuniy loyihani kiritishidan oldin tuzatish bilan hal etiladi.

\section*{E’tirozlarga dalil asosida javob}
Ushbu loyiha uchun 2027-yilgi rasmiy ballar spetsifikatsiyasi olinmagani sababli yangi model aniqligi yoki fiskal ta’siri hisoblanmagan. Lekin ochiq formulalar auditi va kelajakdagi tasdiqlangan hisoblashni tekshirish majburiyati uchun bu fakt to‘siq emas. Manfiy yer misoli operatsion kadastrda uchramasa, bu ma’lumot passportida tekshiriladigan soha bilan asoslanadi; haqiqiy uchrash chastotasi da’vo qilinmagan. Og‘ir vaziyat yo‘li amalda maxsus qoida ustuvorligi bilan hal qilinsa, uning huquqiy asosi va sinov holati qayd etiladi. Maxfiylik shaxsiy ma’lumotlarni tashqi muallifga yuborish orqali emas, qonuniy ichki hisoblash va oshkor etish tekshiruvi orqali ta’minlanadi. Resurs e’tirozi haqiqiy xarajat hisobi bilan ko‘riladi. Mavjud norma ayni talabni qoplasa, ishlab chiquvchi aniq bandni ko‘rsatib, takrorlanadigan normani birlashtirishi mumkin.

\section*{Protsessual kafolatlar va tasdiqlash chegaralari}
\textbf{Mavjud vakolat va zarurat.} Huquqiy qamrov va bajarilgan nazorat alohida tekshiriladi. O‘RQ-682 20-moddasi mavjud qonuniy mexanizm va ma’muriy vakolat masalani hal qilganda unda ko‘rsatilgan loyihalarni tashabbus qilishni cheklaydi; PF-193 yangi baholashning normativ tayyorlanishini bevosita topshiradi. Noaniq huquqiy shartni sinov bayonnomasi yopmaydi. Teng mavjud qoida saqlanadi, texnik integratsiya mavjud vakolatni ishlatadi. 1-materialning besh qisqa bandi yoki teng tahrir to‘liq ilovasiz normativ natijani beradi. Muayyan muallif bandi o‘z-o‘zidan majburiy deb olinmaydi.

\textbf{Murojaatni qabul qilish.} \href{https://lex.uz/docs/-6445145}{Konstitutsiya}ning 40-moddasi vakolatli organlarga taklif berish huquqini va qonuniy tartib hamda muddatlarda ko‘rib chiqishni belgilaydi. \href{https://lex.uz/docs/-3336169}{Murojaatlar haqidagi Qonun}ning 6, 18 va 23-moddalari rekvizitlar, qabul qilish va ro‘yxatdan o‘tkazishni ajratadi. Murojaat kelib tushgan kuni, ish vaqti tugagandan keyin kelsa keyingi ish kuni ro‘yxatdan o‘tkaziladi; ro‘yxatdan o‘tkazishni rad etish taqiqlanadi. Asosiy murojaatda muallif, yashash manzili, oluvchi va masalalar ko‘rsatilgan. 29–30-moddalarda anonimlik va boshqa qonuniy talablarga rioya qilmaslik, qaytarib olish hamda takroriy murojaatning belgilangan shartlari kabi ko‘rib chiqmasdan qoldirish yoki ko‘rib chiqishni tugatish asoslari saqlanadi. Kelajakdagi yuborish yoki avvalgi tashqi yozishmalar haqidagi noma’lum holatlar ushbu tekshiruv bilan tasdiqlanmaydi.

\textbf{Ko‘rib chiqish va javob.} O‘RQ-682 23-moddasi tegishli tadqiqot va murojaatlarni o‘rganish, taklif hamda ilmiy tavsiyalarni umumlashtirish va foydalanishni talab qiladi. 22-modda ishtirokni ruxsat etadi, lekin tayinlashni majburiy qilmaydi. Murojaatlar haqidagi Qonunning 27-moddasi bo‘yicha barcha masalalar ko‘rib chiqilib, yozma yoki elektron javob aniq asoslarni o‘z ichiga olishi kerak. Bandma-band javobni so‘rash shu majburiyatni tekshiriladigan qiladi, ammo qonunda majburiy jadval shakli belgilanmagan. Murojaatlar haqidagi qonun vakolatdan tashqari murojaatni besh kun ichida tegishli organga yuborish va xabar berishni talab etib, asossiz yo‘naltirishni taqiqlaydi. 28-modda taklif uchun bir oygacha muddatni belgilaydi; qo‘shimcha o‘rganishda o‘n kun ichida yozma xabar beriladi. Shaxsiy ishtirok yoki kirish iltimosi alohida murojaat turi va uning muddatini talab qilishi mumkin. Birlashtirilgan xat bu tasnif bo‘yicha majburiy sud talqinini belgilamaydi; qo‘shimcha o‘rganish istisnosiga o‘ylab topilgan yakuniy muddat berilmagan.

\textbf{Bajarilganlik dalili.} Yuborilgan tahrir, ilovalar, qabul yoki elektron tasdiq, ro‘yxat raqami va kelib tushgan sana saqlanadi. Yo‘naltirish va qo‘shimcha o‘rganish xabarlari, qo‘shimcha axborot so‘rovi hamda to‘liq javob saqlanadi. Javob beshta raqamlangan iltimos va har mazmuniy modul bilan solishtiriladi: kiritish, o‘zgartirilgan yechim, aniq amaldagi norma yoki qabul qilmaslik sababi. Manba fayli xeshi qaysi fayl tekshirilganini ko‘rsatadi; organ olganligini isbotlamaydi. Qonuniy qabul qilingan hujjat, rasmiy tayinlash qarori tegishli natijaning dalilidir. Loyihalar kompilyatsiyasi yoki dastur sinovlari bu natijalarni o‘z-o‘zidan hosil qilmaydi.

\textbf{Noqonuniy rad etishdan himoya.} Murojaatlar haqidagi Qonunning 32–33-moddalari qabul qilish yoki ko‘rib chiqishning noqonuniy rad etilishi ustidan shikoyat va shikoyat tartibini tushuntirishni nazarda tutadi; \href{https://lex.uz/docs/-6445145}{Konstitutsiya 55-moddasi} noqonuniy qaror, harakat va harakatsizlik ustidan sud himoyasini ta’minlaydi. \href{https://lex.uz/docs/-3527353}{Ma’muriy sud ishlarini yuritish to‘g‘risidagi kodeks}ning 185–189-moddalari tegishli tartibni belgilaydi. 189-modda bo‘yicha ko‘rsatilgan himoya chorasi uchun noqonuniylik hamda arizachining huquqi yoki qonuniy manfaatiga putur yetishi birga talab qilinadi; sud qonuniy qaror, harakat yoki boshqa bartaraf etish chorasi majburiyatini yuklashi mumkin. Asoslantirilgan siyosiy yechim bilan kelishmaslikning o‘zi huquq buzilishi isboti emas. 186-modda boshqa qonuniy muddat bo‘lmasa, huquq buzilishi ma’lum bo‘lgan paytdan sudga murojaat qilish uchun umumiy olti oylik muddatni belgilaydi; tiklash uzrli sababni talab qiladi. Bu murojaatlar haqidagi Qonunning 22-moddasidagi yuqori turuvchi organga shikoyatning bir yillik muddatidan farq qiladi. Haqiqiy shikoyat paytida faktlar, tasnif, vakolatli sud va amaldagi qoidalar qayta tekshiriladi. Ushbu taklifni qabul qilishga majburlovchi sud qarori olinmagan yoki aniqlanmagan.

\textbf{Huquqiy xulosa doirasi.} O‘RQ-682 23-moddasi muhokama takliflarini tavsiyaviy deb belgilab, ko‘rib chiqishni talab qiladi; 24-modda rad etishga sababini asoslash sharti bilan bevosita ruxsat beradi. Bu yerda dalil bilan asoslanadigan huquqiy kafolat — belgilangan protsessual majburiyatlar, ularning qonuniy shartlari va himoya yo‘llaridir. U tayinlash, cheklangan ma’lumotga kirish, normalarni kiritish yoki qabul qilish kafolati emas. \href{https://lex.uz/docs/-3492199}{O‘RQ-457} 3-moddasi normativ hujjatlarni tayyorlash va qabul qilishni o‘z doirasidan chiqaradi. Shu sababli to‘plamdagi individual ma’muriy qaror talablari hujjatni qabul qildirishning protsessual asosi deb ko‘rsatilmaydi; taklif etilgan murojaat yo‘li fuqaroning murojaati va yuqorida bayon qilingan alohida konstitutsiyaviy hamda qonuniy majburiyatlarga tayanadi.

\textbf{Tatbiq etiladigan tartib va sudda isbotlash.} Murojaatlar haqidagi Qonunning 1-moddasi boshqa qonunchilikda ko‘rib chiqish tartibi belgilangan murojaatlarni o‘z doirasidan chiqaradi. Ushbu xat fuqaroning tashabbuskor taklifi va shaxsiy iltimosi sifatida tayyorlangan, yakunlangan rasmiy muhokama murojaati emas. Organ maxsus tartibni qo‘llasa, so‘ralgan javobda shu tartib va huquqiy asos ko‘rsatilishi lozim; odatiy taklif muddati keyinchalik O‘RQ-682 24-moddasi bo‘yicha portalda beriladigan fikrga tatbiq etilgan deb da’vo qilinmaydi. Ushbu aralash xatning tasnifini belgilovchi majburiy qaror aniqlanmagan. Bu talqin chegarasidir, muallifning konstitutsiyaviy taklif berish huquqi yo‘qligi dalili emas.

Tegishli ma’muriy sud tartibidagi nizoda \href{https://lex.uz/docs/-3527353}{Kodeksning 67-moddasi} isbotlash majburiyatini qarori yoki harakati nizolashilayotgan organ yoki mansabdor shaxsga yuklaydi; e’tiroz asosidagi faktlar tasdiqlanadi va faqat qaror yoki harakatda tayanilgan asoslar keltiriladi. Arizachi o‘z imkoniyati doirasida dalil yig‘ishda qatnashadi va talab qilinayotgan zarar miqdorini isbotlaydi. \href{https://lex.uz/docs/-4711311}{Oliy sud Plenumining 24-son qarori 17 va 20-bandlari} qonuniylikni isbotlash yukini tushuntiradi hamda qonuniy ravishda ma’muriy ixtiyorga berilgan qaror yoki harakatning maqsadga muvofiqligini sud baholamasligini belgilaydi. 21-band belgilangan muddatda ko‘rilmagan yozma murojaat uchun himoya chorasini aniq tushuntiradi: sud harakatsizlikni noqonuniy deb topadi va qanday qaror qabul qilinishini oldindan hal etmasdan, sud belgilagan muddatda ko‘rib chiqish majburiyatini yuklaydi. Shu qoidalar rad etishning haqiqiy asoslari va tegishli majburiyat bajarilishini tekshirishga xizmat qiladi; avtomatik yutuq yoki muallif siyosatini sud tanlashi majburiyatini bermaydi.

\section*{Huquqiy manbalar va yakunlash shartlari}
Mandat: \href{https://lex.uz/docs/-8472526}{PF-193, 13–14-bandlar}. Mavjud mexanizmlar va audit formulalari: \href{https://lex.uz/docs/-8027513}{35-son qaror va ilovalari}. Toifalarning yuqori darajadagi asosi: \href{https://lex.uz/docs/-7952180}{PF-258, 3-band}. Auditdagi \texttt{M}: \href{https://lex.uz/docs/-8283656}{PF-115}; \texttt{P}: \href{https://stat.uz/uz/matbuot-markazi/qo-mita-yangiliklar/67179-minimal-istemol-kharazhatlari-ijmati-t-risida-2}{Milliy statistika qo‘mitasi, 2026-yil 4-mart}. Tayyorlashda ishtirok: \href{https://lex.uz/docs/-5378966}{O‘RQ-682, 22-modda}. Ma’lumotlar va avtomatik qarorlar: \href{https://lex.uz/docs/-4396419}{O‘RQ-547, 11, 16–18, 24 va 28-moddalar; maxsus ma’lumotlar himoyasi}. Individual qaror asoslari: \href{https://lex.uz/docs/-3492199}{O‘RQ-457, 54-modda}. Sezgirlik tahlili uchun metodik manba: \href{https://www.oecd.org/content/dam/oecd/en/publications/reports/2008/08/handbook-on-constructing-composite-indicators-methodology-and-user-guide_g1gh9301/9789264043466-en.pdf}{OECD/JRC qo‘llanmasi, 2008, 7-bosqich}.

Qo‘shimcha manbalar: \href{https://lex.uz/docs/5331933}{PQ-5025, 2-band (b)}, \href{https://lex.uz/docs/-5378966}{O‘RQ-682 ilovasi, 11, 43–51 va 83–85-bandlar}. Analitik matn 11–14 o‘lchamli shrift, normativ bandlar esa 14 o‘lchamli asosiy shrift bilan, A4, chap 3 sm, o‘ng 2 sm va yuqori-past 2 sm hoshiya bilan tayyorlangan. Liberation Serif Times New Roman o‘rnidagi ochiq ko‘rsatilgan almashtirishdir; rasmiy elektron shaklni ishlab chiquvchi tasdiqlashi kerak. Rasmiy LaTeX sinfi topilgani da’vo qilinmaydi.

Bu matn mualliflik taklifidir: yangi majburiyatlar amaldagi huquq sifatida ko‘rsatilmaydi. Yakuniy tahrir uchun rasmiy uslubiyot, barcha tanlangan o‘zgarishlarning yagona rasmiy jadvali, xarajat hisobi, tatbiq etish rekvizitlari hamda o‘zbek huquqiy tili bo‘yicha mutaxassis tekshiruvi zarur. Ushbu modul butun PF-193 topshirig‘ini, xususan xizmatlarga yo‘naltirishning barcha qoidalarini, yakunlamaydi. Qabul qilinishi yoki hech qanday yuridik e’tiroz bo‘lmasligi kafolatlanmaydi.

\end{document}
